\chapter{控制器、主板与反馈：CPU 如何控制整台机器}

前面几章把数据通路讲完了：最小 CPU 章给出了取指、译码、执行的五段路径，主存与总线章给出了地址、数据、握手三类总线事务，存储器件结构章给出了从 SRAM 单元到 DRAM 命令流的完整结构。但一台整机里还有第二张网：CPU 写出一个物理地址，磁盘控制器就开始寻道；设备拉出一条中断线，CPU 就放下当前指令转入服务程序；网卡收到一个包，不经任何一条 CPU 指令就出现在内存里。这些"什么时候让谁动"的决定，构成控制面。值得注意的是，CPU 与外界的一切交互只归结为三类信号：对物理地址的读与写（取指、访存、设备寄存器读写都是同一类事务）、从外部进入的中断输入、设备与内存之间的 DMA 搬运。本章把这张控制网讲透：先看任何设备共有的寄存器接口与状态机，再看 CPU 如何控制自己最快的从设备——内存，然后沿中断链路上行到固件与操作系统，最后把这些分散的机制统一成同一个概念——反馈环。

\begin{keyidea}
控制 = 状态机 + 反馈环。每个控制器（设备、内存控制器、中断控制器）都是一台状态机：它有一组可被软件读写的寄存器作为输入输出接口，有一组时序约束决定状态迁移何时合法。每个控制器又都处在至少一个反馈环里：发出动作、测量结果、按偏差调整。反馈环按响应时间分层——纳秒级的快环（DRAM 时序计数器、PHY 采样点、链路时钟恢复 CDR）必须由硬件闭合，因为软件根本来不及参与；微秒到秒级的慢环（驱动程序的设备管理、固件的内存训练、操作系统的调度）才交给软件。划分的原则只有一条：环的响应时间必须短于被控对象变化的时间常数。
\end{keyidea}

\section{〇、控制面：数据通路之外的第二张网}

最小 CPU 章的数据通路图里，每一条线都有明确的方向和内容：PC 到指令存储器是地址，指令存储器到译码器是指令编码，ALU 到寄存器堆是结果。数据通路回答"数据怎么流动"，但它不回答三个更外层的问题：CPU 怎么让一台设备开始工作？设备怎么让 CPU 知道自己做完了？大批量数据怎么在设备与内存之间移动而不占用流水线？这三个问题的答案都不在数据通路里，而在控制面里。

控制面建立在同一个物理地址空间之上。主存与总线章已经说明，MMIO（Memory-Mapped I/O，内存映射 I/O）把设备寄存器映射进物理地址空间，CPU 用普通的 load/store 就能访问设备；本节要做的是把视角反转过来：从设备一侧看，这些读写不是"访问内存"，而是命令与状态采样。三类信号各管一件事——对物理地址的读与写是一类：MMIO 写是 CPU 发出的命令，MMIO 读是 CPU 对设备状态的采样（取指与访存也属这类事务）；中断是设备反方向发来的通知；DMA 是数据绕开 CPU 流水线的搬运通道。

\begin{pdfdiagram}[x=1cm,y=1cm]
  % CPU 核心居中
  \node[hw compute, minimum width=2.4cm, minimum height=1.4cm] (cpu) at (5.5,0) {CPU 核心\\（流水线/Cache）};
  % 四颗卫星
  \node[hw memory, minimum width=2.6cm, minimum height=1.2cm] (mc) at (1.6,2.6) {内存控制器\\+ DRAM};
  \node[hw control, minimum width=2.6cm, minimum height=1.2cm] (apic) at (9.4,2.6) {中断控制器\\（LAPIC/IOAPIC）};
  \node[hw state, minimum width=2.6cm, minimum height=1.2cm] (pch) at (1.6,-2.6) {PCIe 根复合体\\/ PCH};
  \node[hw state, minimum width=2.6cm, minimum height=1.2cm] (dev) at (9.4,-2.6) {设备 / DMA\\（网卡/磁盘/UART）};
  % 内存读写: CPU -> MC
  \draw[hw bus] (4.3,0.35) -- (3.4,0.35) -- (3.4,2.6) -- (2.9,2.6);
  \node[hw label, font=\tiny, anchor=west, text width=2.4cm] at (3.55,1.5) {内存读写\\（物理地址）};
  % 中断: APIC -> CPU
  \draw[hw return] (8.1,2.6) -- (7.6,2.6) -- (7.6,0.35) -- (6.7,0.35);
  \node[hw label, font=\tiny, anchor=west] at (7.75,1.5) {中断请求};
  % MMIO: CPU -> PCH
  \draw[hw bus] (4.3,-0.35) -- (3.2,-0.35) -- (3.2,-2.6) -- (2.9,-2.6);
  \node[hw label, font=\tiny, anchor=west, text width=2.4cm] at (3.35,-1.5) {MMIO 读写\\（设备寄存器）};
  % 中断: 设备 -> APIC
  \draw[hw return] (9.4,-2.0) -- (9.4,2.0);
  \node[hw label, font=\tiny, anchor=east] at (9.25,0) {中断（MSI/引脚）};
  % DMA 数据: 设备 -> PCH -> MC
  \draw[hw bus] (9.4,-3.2) -- (9.4,-3.55) -- (1.6,-3.55) -- (1.6,-3.2);
  \node[hw label, font=\tiny] at (5.5,-3.85) {DMA 数据（直达内存，不占用 CPU 流水线）};
  \draw[hw bus] (1.6,-2.0) -- (1.6,2.0);
  \node[hw label, font=\tiny, anchor=west] at (1.75,0) {DMA 数据};
\end{pdfdiagram}
\diagramnote{控制面总图。CPU 核心居中，向外辐射三类线：实线是地址与数据路径——CPU 对 DRAM 的内存读写、对设备寄存器的 MMIO 读写、设备到内存的 DMA 数据；虚线是反方向的事件通知——设备的中断请求经中断控制器汇总后投递给 CPU。DMA 数据经根复合体直达内存控制器，途中不进入任何一条 CPU 指令的执行路径。}

\topic{写寄存器是命令，读寄存器是状态}

MMIO 写与普通内存写有本质区别。往 DRAM 地址写一个值，语义是"保存这个值，以后读回来"；往设备控制寄存器写一个值，语义是"按这个值行动"——写 \code{0x01} 可能表示启动一次发送，写 \code{0x02} 可能表示复位状态机。写的值不一定要被保存，写的动作本身才是信息。主存与总线章已经指出，MMIO 写的总线响应只说明命令到达设备接口，不代表设备动作完成；从控制面看，这句话可以换个说法：MMIO 写是一条单向命令通道，完成与否要用另一条通道（状态读或中断）来确认。

MMIO 读则是对设备状态的一次采样。状态寄存器由设备硬件持续更新，CPU 读到的只是采样时刻的快照。两次读之间状态可能任意变化，这与普通内存"读到的就是上次写的"完全不同。第一节会展开这组差异如何塑造寄存器接口的各种读写约定。

\topic{中断是反方向的回信}

命令与状态两个通道都是 CPU 主动的：CPU 决定何时写、何时读。但设备完成一件事的时间点是设备才知道的，若只能靠 CPU 反复读状态去发现，CPU 就被拴在慢设备上。中断把方向反转：设备在状态变化时主动向 CPU 发一个事件信号，中断控制器负责汇总、屏蔽和排序，CPU 在指令边界接收。主存与总线章第四节已经走完"事件→状态位→IRQ→控制器→向量→服务程序"的完整链路；本章第三节会把中断控制器本身作为控制面的一台核心状态机拆开。眼下只需记住总图里的方向约定：实线由 CPU 或 DMA 发起，虚线由设备发起，控制面是双向的。

\section{一、控制的基本元件：寄存器接口与状态机}

\topic{1.1\quad 任何设备 = 一组寄存器 + 一台状态机}

把网卡、磁盘控制器、UART、定时器、GPU 命令处理器切开，结构惊人地一致：面向 CPU 的是一组寄存器，寄存器后面是一台状态机，状态机驱动设备的引脚或内部数据通路。寄存器按角色分三类：

- \hwkey{控制寄存器}（Control Register, CR）：软件写入，作为状态机的输入。写值告诉状态机做什么——启动、停止、选模式、开中断。
- \hwkey{状态寄存器}（Status Register, SR）：硬件写入、软件读出，是状态机对外的输出。读值告诉软件状态机现在在哪——忙不忙、有没有错、数据到了没有。
- \hwkey{数据寄存器}（Data Register, DR）：数据通路的出入口。写它把数据推入发送 FIFO，读它从接收 FIFO 弹出数据。数据寄存器通常有副作用：读一次少一个元素，写一次多一个元素。

这组寄存器怎样被物理地址选中？答案与最小 CPU 章的指令译码器是同一种电路：那里是 opcode 进译码器、控制线出；这里是物理地址高位进译码器、片选出。地址高位与设备被分配的地址区间比较，命中则产生片选使能该设备的寄存器堆，地址低位直接作为寄存器号在堆内选择——经典芯片与平台章提到的 74LS138 三八译码器就是这套机制在分立元件时代的形态，今天它集成在根复合体和每个设备的地址译码逻辑里。软件一侧看到的对应物是驱动的 \code{ioremap}/\code{readl}/\code{writel}：\code{ioremap} 把这段物理区间映射进内核虚拟地址，并把它建立为不可缓存（UC）的映射（属性由 PAT/MTRR 保证）；\code{readl}/\code{writel} 则提供 volatile 语义与必要的顺序、屏障保证，使访问真正发到设备。

\begin{pdfdiagram}[x=1cm,y=1cm]
  % 地址译码器
  \node[hw compute, minimum width=2.4cm, minimum height=1.0cm] (dec) at (1.7,2.8) {地址译码器\\（区间比较）};
  \draw[hw bus] (-0.7,2.8) -- (0.5,2.8);
  \node[hw label, font=\tiny, anchor=south] at (-0.35,2.86) {地址高位};
  % 寄存器堆
  \node[hw state, minimum width=3.2cm, minimum height=2.6cm] (rf) at (5.0,0.6) {};
  \draw[BookTeal, line width=0.4pt] (3.7,1.25) rectangle (6.3,1.75);
  \node[hw label, font=\tiny] at (5.0,1.5) {控制寄存器 CR};
  \draw[BookTeal, line width=0.4pt] (3.7,0.45) rectangle (6.3,0.95);
  \node[hw label, font=\tiny] at (5.0,0.7) {状态寄存器 SR};
  \draw[BookTeal, line width=0.4pt] (3.7,-0.35) rectangle (6.3,0.15);
  \node[hw label, font=\tiny] at (5.0,-0.1) {数据寄存器 / FIFO};
  % 片选
  \draw[hw bus] (2.9,2.8) -- (3.8,2.8) -- (3.8,1.9);
  \node[hw label, font=\tiny, anchor=south] at (3.35,2.9) {片选};
  % 低位地址 / 数据
  \draw[hw bus] (-0.7,1.2) -- (3.4,1.2);
  \node[hw label, font=\tiny, anchor=south] at (0.5,1.28) {地址低位（寄存器号）};
  \draw[hw bus] (-0.7,0.5) -- (3.4,0.5);
  \node[hw label, font=\tiny, anchor=south] at (0.5,0.58) {写数据};
  \draw[hw bus] (3.4,-0.1) -- (-0.7,-0.1);
  \node[hw label, font=\tiny, anchor=north] at (0.5,-0.18) {读数据};
  % 状态机
  \node[hw control, minimum width=2.6cm, minimum height=2.6cm] (fsm) at (9.3,0.6) {设备状态机};
  \draw[hw bus] (6.6,1.5) -- (8.0,1.5);
  \node[hw label, font=\tiny, anchor=south] at (7.3,1.56) {控制位};
  \draw[hw return] (8.0,0.7) -- (6.6,0.7);
  \node[hw label, font=\tiny, anchor=south] at (7.3,0.76) {状态位};
  \draw[hw bus] (6.6,-0.1) -- (8.0,-0.1);
  \node[hw label, font=\tiny, anchor=south] at (7.3,-0.04) {发送数据};
  \draw[hw bus] (8.0,-0.45) -- (6.6,-0.45);
  \node[hw label, font=\tiny, anchor=north] at (7.3,-0.51) {接收数据};
  % 引脚
  \draw[hw bus] (10.6,0.9) -- (11.4,0.9);
  \node[hw label, font=\tiny, anchor=south] at (11.0,0.96) {输出引脚};
  \draw[hw bus] (11.4,0.25) -- (10.6,0.25);
  \node[hw label, font=\tiny, anchor=north] at (11.0,0.19) {输入引脚};
  % IRQ
  \draw[hw return] (9.3,1.9) -- (9.3,2.9);
  \node[hw label, font=\tiny, anchor=west] at (9.45,2.4) {IRQ};
\end{pdfdiagram}
\diagramnote{设备寄存器接口。地址高位经译码器产生片选选中整台设备，地址低位在寄存器堆内选择寄存器，与最小 CPU 章的 opcode 译码是同一种电路。CR 是状态机的输入，SR 是状态机的输出，数据寄存器/FIFO 连接状态机内部的数据通路，IRQ 是状态机向 CPU 方向的唯一主动信号。}

\topic{1.2\quad 寄存器接口的读写约定：为什么会有这些奇怪规则}

翻开任何一份设备数据手册，寄存器位描述里总有一列读写属性，其中不少看起来不合直觉：有的位"读了就清零"，有的位"写 1 清零、写 0 无效"，有的位"置位后由硬件自己清回 0"。这些约定不是设计者的癖好，全部来自同一个结构性事实：\hwkey{寄存器有两个并发写者}——软件和硬件。普通内存只有一个写者（程序），设备寄存器却是软件与设备状态机共享的变量，竞态随时可能发生。

\begin{longtable}{L{0.21\linewidth}L{0.36\linewidth}L{0.36\linewidth}}
\toprule
约定 & 行为 & 为什么存在 \\
\midrule
read-to-clear & 读该寄存器后某些位自动清零 & 中断事件位：读走事件即表示"已取走"，硬件随即清位准备记录下一事件；免去了软件再写一次清除 \\
W1C（写 1 清零） & write-1-to-clear：写 1 的位清零，写 0 的位不变 & 让软件只清自己处理的位，不清并发到达的新事件位 \\
自清零（self-clearing） & 软件写 1 启动动作，硬件在动作开始后自动清回 0 & 启动位表示"请求"而非"状态"：软件读回 0 即知硬件已接手，避免软件忘清导致重复启动 \\
只读（RO）/ 只写（WO） & 单向寄存器 & 状态位被软件写会制造假象；命令位被软件读回没有物理来源——索性在硬件上禁止 \\
\bottomrule
\end{longtable}

W1C 值得展开，因为它直接回应竞态。典型错误是"读-改-写"：软件读出 \code{0x05}（位 0 和位 2 有事件），处理完位 0 后清掉该位得到 \code{0x04}，再写回——若在读与写之间硬件刚置起了位 1，这次写回会把新事件位又抹回 0，事件被无声丢失。W1C 从根上消除这个窗口：写 \code{0x01} 只清位 0，写 0 的位硬件不理会，于是软件的清除动作与硬件的置位动作在位一级互不干涉，不需要任何锁。

\begin{warningbox}
把设备状态寄存器当普通变量做"读-改-写"是驱动开发中反复出现的缺陷来源。C 里的 \code{*p |= FLAG} 编译出来正是读-改-写三条指令，对 W1C 寄存器执行 \code{|=} 会把读到的所有事件位全部清掉（每个读到的 1 都被写回）。正确做法是直接把要清的位作为常量写入：\code{*p = FLAG}。另外，\code{volatile} 只阻止编译器优化掉这次访问，它管不了硬件并发——约定是硬件给的，遵守靠软件。
\end{warningbox}

\topic{1.3\quad 轮询还是中断：一个代价模型}

设备完成一件事之后，软件有两种发现方式：反复读状态寄存器（轮询），或等设备发中断。两者各有确定的代价，选择是一个定量问题。

轮询的代价是持续的 CPU 占用。每次查询是一次 MMIO 读：读不能走 cache，必须发到设备再返回，PCIe 设备上一次 MMIO 读的往返在几百纳秒到 1 微秒量级。若每 $P$ 秒查询一次，CPU 在查询上的占用率约为 $T_q/P$（$T_q$ 是单次查询耗时），而事件被发现的最大延迟是 $P$。$P$ 取得小，发现得快但占用率线性上升；$P$ 取得大，省 CPU 但响应慢。中断的代价是每次事件一笔固定开销：CPU 保存现场、经向量表分发、进入服务程序，再恢复现场返回，全程在微秒量级，事件稀疏时几乎为零成本，但事件密集到接近这笔开销的倒数时，CPU 会被中断进入/退出本身耗尽——这正是网络设备在中断风暴下吞吐反而下降的原因。

\begin{longtable}{L{0.16\linewidth}L{0.38\linewidth}L{0.38\linewidth}}
\toprule
维度 & 轮询 & 中断 \\
\midrule
CPU 开销 & 与事件无关，$T_q/P$ 持续占用 & 与事件率成正比，每次约微秒量级 \\
发现延迟 & 最坏一个查询周期 $P$，可预期 & 进入延迟加调度抖动，波动较大 \\
密集事件 & 开销不随事件增长，反而划算 & 每次事件一笔固定开销，可能耗尽 CPU \\
稀疏事件 & 绝大多数查询空手而归 & 几乎零成本 \\
\bottomrule
\end{longtable}

工程上的选择规则因此是：事件间隔远小于中断开销时轮询（高性能存储、用户态网络栈的 busy-poll），事件稀疏且到达时间不可预知时中断（键盘、串口），突发型负载用混合策略——Linux 网卡的 NAPI 是教科书式的例子：第一个包触发中断，服务程序关掉中断改为轮询，把积压的包一批收完，队列空了再重新开中断。中断负责"叫醒"，轮询负责"趁醒着多干"，两种机制的切换点由队列是否腾空这个状态决定——这本身就是一个小小的反馈环。

\topic{1.4\quad 两个经典例子：8254 PIT 与 16550 UART}

\textbf{8254 可编程间隔定时器}（Programmable Interval Timer, PIT）把前面所有概念集于一身。它内含三个独立的 16 位递减计数器，共享一个 \hwevidence{1.193182 MHz} 的输入时钟；每个计数器装入初值 $N$ 后逐拍递减，数到 0 时按所选方式动作，并在方式 2/3 下自动重装初值。方式 3（方波发生器）下输出引脚 OUT 产生周期为 $N$ 个输入时钟的方波——分频器就是这么来的。IBM PC 把计数器 0 的输出接到 8259 的 IR0，成为操作系统的时钟节拍来源：初值 0 按 65536 计，$1193182/65536 \approx 18.2$ Hz 的节拍是 PC 兼容机沿用多年的传统。

\begin{longtable}{L{0.14\linewidth}L{0.24\linewidth}L{0.55\linewidth}}
\toprule
端口 & 寄存器 & 作用 \\
\midrule
\code{0x40} & 计数器 0 & 读：当前计数值；写：初值低/高字节 \\
\code{0x41} & 计数器 1 & 同上（PC 上曾用于 DRAM 刷新节拍） \\
\code{0x42} & 计数器 2 & 同上（PC 上接扬声器） \\
\code{0x43} & 控制字寄存器（WO） & 选择计数器、读写顺序、工作方式、数制 \\
\bottomrule
\end{longtable}

控制字（写 \code{0x43}）的高两位选计数器，中间两位选读写顺序（先低字节后高字节；00 为锁存命令，即下文的计数值快照），接下来三位选方式 0--5，最低位选二进制或 BCD（Binary-Coded Decimal，二-十进制编码）计数。把计数器 0 配成方式 3、输出约 1 kHz 方波的完整序列只有六条指令（x86 端口 I/O，\code{out dx, al} 中端口地址经 \code{dx} 给出）：

\begin{lstlisting}
mov al, 0x36        ; 计数器0, 先低后高, 方式3, 二进制
out 0x43, al
mov al, 0xA9        ; 初值低字节
out 0x40, al
mov al, 0x04        ; 初值高字节: 0x04A9 = 1193
out 0x40, al        ; 1193182 / 1193 ≈ 1000 Hz
\end{lstlisting}

8254 还有一个细节恰好印证 1.2 节的并发问题：计数器在不停地减，软件读 16 位计数值要分两次读，两次读之间低字节可能向高字节借位，拼出来的值就是错的。解法是锁存命令——先向 \code{0x43} 发一条锁存命令，硬件把当前计数值瞬时复制到一个保持寄存器，软件再分两字节读这份快照。硬件更新与软件读取的竞态，由硬件提供一个原子快照来化解。

\begin{classicnote}
1.193182 MHz 这个古怪的频率来自 IBM PC 的 14.31818 MHz 晶体——它是 NTSC 彩色副载波 3.579545 MHz 的 4 倍，同一颗晶体既供显示又供 CPU 外设时钟；8254 的输入是它除以 12。经典芯片与平台章讲过这种"一颗晶体供全机"的时代，HPET 出现之前的二十多年里，操作系统的时间观念就建立在这 18.2 Hz 的节拍上。
\end{classicnote}

\textbf{16550 UART} 是串口控制器，把 CPU 的并行字节与串行线上的位流互相转换（经典芯片与平台章已在平台清单里提到它，这里拆寄存器）。它以 8 个端口偏移呈现一组寄存器，多数偏移上读和写是两个不同的寄存器：

\begin{longtable}{L{0.10\linewidth}L{0.10\linewidth}L{0.30\linewidth}L{0.36\linewidth}}
\toprule
偏移 & DLAB & 读 & 写 \\
\midrule
0 & 0 & RBR：接收缓冲（RO，读了就取走） & THR：发送保持（WO） \\
0 & 1 & DLL：除数低字节 & DLL：除数低字节 \\
1 & 0 & IER：中断使能 & IER：中断使能 \\
1 & 1 & DLH：除数高字节 & DLH：除数高字节 \\
2 & -- & IIR：中断标识（RO，读即应答） & FCR：FIFO 控制（WO） \\
3 & -- & LCR：线路控制（位 7 即 DLAB） & LCR：线路控制 \\
4 & -- & MCR：调制解调器控制 & MCR：调制解调器控制 \\
5 & -- & LSR：线路状态（RO） & -- \\
6 & -- & MSR：调制解调器状态（RO） & -- \\
7 & -- & SCR：暂存 & SCR：暂存 \\
\bottomrule
\end{longtable}

这张表本身就是 1.2 节约定的陈列柜：RBR 是 read-to-clear 的数据寄存器，IIR 是 read-to-acknowledge 的状态寄存器，THR 只写、LSR 只读。DLAB（Divisor Latch Access Bit，除数锁存访问位，LCR 的位 7）是另一种约定——偏移 0/1 各有两个目标寄存器，地址位不够用了，就拿一个模式位当"第 9 位地址"：DLAB=1 时访问除数锁存器，DLAB=0 时访问数据与中断使能。波特率由除数决定：输入时钟 1.8432 MHz，先固定除以 16 再除以除数锁存器的值，除数为 1 即 \hwevidence{115200 baud}。初始化 115200 8N1 并打开 FIFO 的序列：

\begin{lstlisting}
; COM1 = 0x3F8
mov dx, 0x3FB       ; LCR
mov al, 0x80        ; DLAB=1
out dx, al
mov dx, 0x3F8       ; DLL
mov al, 1           ; 除数 1 -> 115200 baud
out dx, al
mov dx, 0x3F9       ; DLH = 0
xor al, al
out dx, al
mov dx, 0x3FB       ; LCR: DLAB=0, 8 数据位/无校验/1 停止位
mov al, 0x03
out dx, al
mov dx, 0x3FA       ; FCR: 使能 FIFO, 清收发 FIFO, 触发水位 14
mov al, 0xC7
out dx, al
\end{lstlisting}

发送时的轮询循环直接对应 1.3 节的代价模型——每次 \code{in} 都是一次端口周期，CPU 在 \code{.wait} 上空转的拍数就是查询开销：

\begin{lstlisting}
; bl = 待发送字节
mov dx, 0x3FD       ; LSR
.wait:
in  al, dx          ; 读状态寄存器 = 一次状态采样
test al, 0x20       ; THRE: 发送保持寄存器空?
jz  .wait
mov dx, 0x3F8       ; THR
mov al, bl
out dx, al
\end{lstlisting}

LSR 的关键位：

\begin{longtable}{L{0.10\linewidth}L{0.16\linewidth}L{0.66\linewidth}}
\toprule
位 & 名称 & 含义 \\
\midrule
0 & DR & 接收数据就绪：RBR 里有未读字节（读 RBR 后由硬件清零） \\
1 & OE & 溢出错误：新字节到达时 RBR 还没被读走，旧字节被覆盖 \\
5 & THRE & 发送保持寄存器空：可以接受下一个待发送字节 \\
6 & TEMT & 发送器完全空：移位寄存器也空了，最后一位已真正上线 \\
\bottomrule
\end{longtable}

16550 相对早期 8250 的主要改进是 16 字节收发 FIFO，而 FCR 的触发水位（1/4/8/14 字节可选）体现的是同一个权衡：水位设 14，则收满 14 字节才发一次中断，中断频率降为逐字节中断的十几分之一，但软件看到数据的延迟最长可达 14 个字节时间。它就是 1.3 节"查询周期 $P$"在硬件里的翻版，只是这个环由硬件闭合，软件只在配置时选一次参数。

\section{二、CPU 如何控制内存：地址译码、命令生成与训练}

内存是 CPU 最特殊的控制对象：它不是外设，却比任何外设都更需要控制。DRAM 芯片自己什么都不记得——哪一行开着、隔多久才能发下一条命令、采样点定在哪个相位，全部由 CPU 一侧的内存控制器维护。本节按控制的三个层次展开：先控制"去哪"（物理地址地图），再控制"何时动"（命令生成与调度），最后控制"参数取多少"（上电训练）。

\topic{2.1\quad 物理地址空间地图：DRAM 窗口、MMIO 洞与固件区}

物理地址不是 DRAM 的同义词。一个 32 位物理地址送出 CPU 后，第一级译码要回答的问题是：这个地址属于 DRAM、某台设备的寄存器窗口，还是固件 ROM？主存与总线章已经给出 MMIO 不可缓存等属性要求，经典芯片与平台章给出 PCIe 枚举时如何把各设备的 BAR（Base Address Register，基址寄存器）窗口分配到互不重叠的区间；这里把结果画成一张完整地图。

\begin{pdfdiagram}[x=1cm,y=1cm]
  % 主条：x 0..1.8, y 0..10.2（非线性示意）
  % 常规内存 0-640KB
  \fill[BookTeal!14] (0,0) rectangle (1.8,1.0);
  \draw[BookRule, line width=0.5pt] (0,0) rectangle (1.8,1.0);
  % VGA 0xA0000-0xBFFFF
  \fill[BookAmber!18] (0,1.0) rectangle (1.8,1.6);
  \draw[BookRule, line width=0.5pt] (0,1.0) rectangle (1.8,1.6);
  % PAM 区 0xC0000-0xFFFFF
  \fill[BookRed!10] (0,1.6) rectangle (1.8,2.4);
  \draw[BookRule, line width=0.5pt] (0,1.6) rectangle (1.8,2.4);
  % 主 DRAM 1MB-TOLUD
  \fill[BookTeal!14] (0,2.4) rectangle (1.8,5.6);
  \draw[BookRule, line width=0.5pt] (0,2.4) rectangle (1.8,5.6);
  % PCIe MMIO 洞 TOLUD-0xEFFFFFFF
  \fill[BookAmber!18] (0,5.6) rectangle (1.8,7.6);
  \draw[BookRule, line width=0.5pt] (0,5.6) rectangle (1.8,7.6);
  % HPET/IOAPIC/LAPIC
  \fill[BookAccent!12] (0,7.6) rectangle (1.8,8.3);
  \draw[BookRule, line width=0.5pt] (0,7.6) rectangle (1.8,8.3);
  % 固件 ROM
  \fill[BookRed!10] (0,8.3) rectangle (1.8,9.2);
  \draw[BookRule, line width=0.5pt] (0,8.3) rectangle (1.8,9.2);
  % 4GB 以上重映射
  \fill[BookTeal!14] (0,9.2) rectangle (1.8,10.2);
  \draw[BookRule, line width=0.5pt] (0,9.2) rectangle (1.8,10.2);
  % 地址轴
  \draw[BookMuted, line width=0.6pt, ->] (-0.5,0) -- (-0.5,10.5);
  \node[hw label, font=\tiny] at (-0.5,10.75) {地址增大};
  % 右侧标注（引导线 + 文字）
  \draw[BookMuted, line width=0.3pt] (1.8,0.5) -- (2.1,0.5);
  \node[hw label, font=\tiny, anchor=west] at (2.2,0.5) {\code{0x00000000}--\code{0x0009FFFF} 常规内存 640 KB（DRAM）};
  \draw[BookMuted, line width=0.3pt] (1.8,1.3) -- (2.1,1.3);
  \node[hw label, font=\tiny, anchor=west] at (2.2,1.3) {\code{0x000A0000}--\code{0x000BFFFF} VGA 帧缓冲（MMIO）};
  \draw[BookMuted, line width=0.3pt] (1.8,2.0) -- (2.1,2.0);
  \node[hw label, font=\tiny, anchor=west] at (2.2,2.0) {\code{0x000C0000}--\code{0x000FFFFF} 扩展 ROM/BIOS 区（PAM 控制）};
  \draw[BookMuted, line width=0.3pt] (1.8,4.0) -- (2.1,4.0);
  \node[hw label, font=\tiny, anchor=west] at (2.2,4.0) {\code{0x00100000}--TOLUD 主 DRAM 窗口（可达数 GB）};
  \draw[BookMuted, line width=0.3pt] (1.8,6.6) -- (2.1,6.6);
  \node[hw label, font=\tiny, anchor=west] at (2.2,6.6) {TOLUD--\code{0xEFFFFFFF} PCIe MMIO 洞（各设备 BAR 窗口）};
  \draw[BookMuted, line width=0.3pt] (1.8,7.95) -- (2.1,7.95);
  \node[hw label, font=\tiny, anchor=west] at (2.2,7.95) {\code{0xFED00000} HPET、\code{0xFEC00000} IOAPIC、\code{0xFEE00000} LAPIC};
  \draw[BookMuted, line width=0.3pt] (1.8,8.75) -- (2.1,8.75);
  \node[hw label, font=\tiny, anchor=west] at (2.2,8.75) {\code{0xFF000000}--\code{0xFFFFFFFF} 固件 ROM（复位向量 \code{0xFFFFFFF0}）};
  \draw[BookMuted, line width=0.3pt] (1.8,9.7) -- (2.1,9.7);
  \node[hw label, font=\tiny, anchor=west] at (2.2,9.7) {4 GB 以上：被 MMIO 洞挡住的 DRAM 重映射到此};
\end{pdfdiagram}
\diagramnote{典型 PC 物理地址地图（非线性比例）。绿色是 DRAM 可达的窗口，黄色是 MMIO，红色是 ROM/固件区；顶部 \code{0xFEC00000}--\code{0xFEE00000} 一带是 IOAPIC、HPET、LAPIC 各控制器的固定 MMIO 窗口区。1 MB 以下的布局是 PC 兼容遗产；TOLUD 以下尽量留给 DRAM，TOLUD 到 4 GB 之间是 MMIO 洞；被洞挡住的物理内存由重映射寄存器搬到 4 GB 以上，64 位系统照常可用。}

这张地图不是硬件写死的，而是固件在上电早期配置的。Intel 平台的两组寄存器决定了地图的形状：\hwkey{TOLUD}（Top of Low Usable DRAM，低位可用 DRAM 顶界）由固件按实际内存容量写入 host bridge，它以下的地址译码到 DRAM、以上（至 4 GB）译码到 PCIe/MMIO；\hwkey{PAM}（Programmable Attribute Map，可编程属性映射）寄存器组把 \code{0xC0000}--\code{0xFFFFF} 这段 256 KB 按 16 KB 粒度逐段设置读写属性——指向 ROM 只读、指向 DRAM 可读写、或谁都不响应。系统 BIOS 区默认指向固件 ROM，固件跑起来之后常把 ROM 内容复制到同地址的 DRAM 并把该段 PAM 切到 DRAM（shadow RAM），因为 DRAM 比固件芯片快得多。地图顶端的几个地址是各控制器的固定窗口：本地 APIC 在 \code{0xFEE00000}、IOAPIC 在 \code{0xFEC00000}、HPET（High Precision Event Timer，高精度事件定时器）在 \code{0xFED00000}，最顶端 16 MB 留给固件，复位后 CPU 取的第一条指令在 \code{0xFFFFFFF0}——第四节讲启动时还会回到这里。

ARM 世界没有这套 PC 遗产，地址地图由 SoC 厂商自行定义：片上 SRAM、DRAM 窗口、各外设寄存器块的位置每颗芯片都不同，内核无法靠猜。于是发现机制也不同——x86 靠固件提供的 E820/ACPI 表，ARM 靠\term{设备树}{device tree}：一个随固件传给内核的数据结构，\code{/memory} 节点声明 DRAM 窗口，各总线节点的 \code{reg} 属性声明每个设备的寄存器基址与长度。地图的形状不同，问题相同：任何一个物理地址都必须有唯一确定的响应者，或一个确定的"无响应"错误——这正是 1.1 节译码器在整机尺度的展开。

\topic{2.2\quad 从"读 64 字节"到命令序列：约束计数器与一拍决策}

存储器件结构章第四节已经把内存控制器的算法面讲完了：bank 状态机、全套时序参数、FR-FCFS 调度策略；第八节的 trace 给出了从 load 指令到 DRAM 命令的端到端路径。本节不重复这些内容，改问一个实现问题：这些算法在硬件里到底长什么样？一条"读 64 字节"的请求进了控制器，调度器凭什么在一拍之内说出"现在能发哪条命令"？

答案的第一部分是\hwkey{约束计数器}。DRAM 的每条时序约束（存储器件结构章第四节那张参数表里的 tRCD、tRAS、tRP、tRRD、tFAW 等）含义都是"某条命令之后，多少拍之内不许发某类命令"。硬件把这个含义直译成电路：每个 bank 挂一列倒计数器，控制器每发出一条命令，就把相关参数值装进目标 bank 对应的计数器——发 ACT 时装入 \hwtiming{tRCD} 与 \hwtiming{tRAS}，发 PRE 时装入 \hwtiming{tRP}；此后每个时钟周期所有计数器同时减一，减到 0 即表示该约束解除。一个 bank 此刻能否接受 RD/WR，看它的 tRCD 计数器是否为 0；能否接受 PRE，看 tRAS 是否到期——每条约束一个比较器，全部比较器的结果与起来就是一个 ready 位。tFAW 这类跨 bank 的全局约束则另有实现：维护一个装最近 4 次 ACT 时间戳的小移位链，最早那次超过 \hwtiming{tFAW}=34 CK 窗口之前不再批准新的 ACT。

答案的第二部分是并行。队列里几十条请求，每条都要算自己的 ready 位（目标 bank 的状态与计数器全在片内寄存器里，读取没有延迟），几十个 ready 位同拍得出；随后是一级优先级编码器——先筛 row-hit，再在命中者中取时间戳最老者，这正是存储器件结构章第四节 FR-FCFS 算法的电路形态：它不是一段顺序执行的程序，而是一拍内稳定下来的组合网络。调度"算法"与"电路"的对应关系，和最小 CPU 章里"指令语义"与"控制字真值表"的关系如出一辙。

被选中的请求最后交给命令编码器，翻译成 CA 总线上的一个周期。DDR 命令用几根控制线的电平组合编码（下表为 DDR3 风格的经典编码，信号均为低有效；DDR4/DDR5 引脚定义不同，如 DDR4 有独立的 ACT\_n，但"组合编码 + 时序约束"的结构相同）：

\begin{longtable}{L{0.13\linewidth}L{0.09\linewidth}L{0.09\linewidth}L{0.09\linewidth}L{0.09\linewidth}L{0.39\linewidth}}
\toprule
命令 & CS & RAS & CAS & WE & 含义 \\
\midrule
NOP & 0 & 1 & 1 & 1 & 空操作，保持总线同步 \\
ACT & 0 & 0 & 1 & 1 & 激活：打开指定 bank 的一行进 row buffer \\
RD & 0 & 1 & 0 & 1 & 读：从打开行取一列，CL 拍后数据上线 \\
WR & 0 & 1 & 0 & 0 & 写：向打开行写一列，WL 拍后发出数据 \\
PRE & 0 & 0 & 1 & 0 & 预充：关闭打开的行，tRP 后回 IDLE \\
REF & 0 & 0 & 0 & 1 & 按内部刷新计数器刷新全部 bank 的对应行，占用 tRFC \\
\bottomrule
\end{longtable}

\begin{pdfdiagram}[x=1cm,y=1cm]
  % bank 状态表：4 行 x 4 列 + 就绪位
  \draw[BookRule, line width=0.5pt] (0,1.1) rectangle (4.6,3.85);
  \draw[BookRule, line width=0.3pt] (0,3.3) -- (4.6,3.3);
  \draw[BookRule, line width=0.3pt] (0,2.75) -- (4.6,2.75);
  \draw[BookRule, line width=0.3pt] (0,2.2) -- (4.6,2.2);
  \draw[BookRule, line width=0.3pt] (0,1.65) -- (4.6,1.65);
  \draw[BookRule, line width=0.3pt] (0.7,1.1) -- (0.7,3.85);
  \draw[BookRule, line width=0.3pt] (1.8,1.1) -- (1.8,3.85);
  \draw[BookRule, line width=0.3pt] (2.8,1.1) -- (2.8,3.85);
  % 表头
  \node[hw label, font=\tiny\bfseries] at (0.35,3.58) {bank};
  \node[hw label, font=\tiny\bfseries] at (1.25,3.58) {状态};
  \node[hw label, font=\tiny\bfseries] at (2.3,3.58) {打开行};
  \node[hw label, font=\tiny\bfseries] at (3.7,3.58) {tRCD/tRAS/tRP\\计数};
  % 行内容
  \node[hw label, font=\tiny] at (0.35,3.03) {0};
  \node[hw label, font=\tiny] at (1.25,3.03) {ACT};
  \node[hw label, font=\tiny] at (2.3,3.03) {\code{0x1F3}};
  \node[hw label, font=\tiny] at (3.7,3.03) {0 / 12 / 0};
  \node[hw label, font=\tiny] at (0.35,2.48) {1};
  \node[hw label, font=\tiny] at (1.25,2.48) {ACT};
  \node[hw label, font=\tiny] at (2.3,2.48) {\code{0x0A2}};
  \node[hw label, font=\tiny] at (3.7,2.48) {0 / 7 / 0};
  \node[hw label, font=\tiny] at (0.35,1.93) {2};
  \node[hw label, font=\tiny] at (1.25,1.93) {IDLE};
  \node[hw label, font=\tiny] at (2.3,1.93) {---};
  \node[hw label, font=\tiny] at (3.7,1.93) {0 / 0 / 0};
  \node[hw label, font=\tiny] at (0.35,1.38) {3};
  \node[hw label, font=\tiny] at (1.25,1.38) {PRE};
  \node[hw label, font=\tiny] at (2.3,1.38) {---};
  \node[hw label, font=\tiny] at (3.7,1.38) {-- / -- / 4};
  % 就绪位小格
  \foreach \y in {3.03, 2.48, 1.93, 1.38} {
    \draw[BookRule, line width=0.3pt] (4.9,\y-0.22) rectangle (5.9,\y+0.22);
  }
  \node[hw label, font=\tiny, text=BookTeal] at (5.4,3.03) {RD 就绪};
  \node[hw label, font=\tiny, text=BookTeal] at (5.4,2.48) {RD 就绪};
  \node[hw label, font=\tiny, text=BookTeal] at (5.4,1.93) {ACT 就绪};
  \node[hw label, font=\tiny, text=BookRed] at (5.4,1.38) {等待};
  % 连线到仲裁器
  \draw[hw bus] (5.9,3.03) -- (6.3,3.03) -- (6.3,2.6) -- (6.6,2.6);
  \draw[hw bus] (5.9,2.48) -- (6.6,2.48);
  \draw[hw bus] (5.9,1.93) -- (6.6,1.93);
  \draw[hw return] (5.9,1.38) -- (6.6,1.38);
  % 仲裁器
  \node[hw control, minimum width=2.4cm, minimum height=1.8cm] (arb) at (7.8,1.9) {优先级仲裁\\row-hit 优先\\同级 FCFS};
  % 命令编码器
  \node[hw compute, minimum width=2.0cm, minimum height=1.8cm] (enc) at (10.7,1.9) {命令编码\\CS/RAS/\\CAS/WE};
  \draw[hw bus] (9.0,1.9) -- (9.7,1.9);
  \node[hw label, font=\tiny, anchor=south] at (9.4,1.94) {选中};
  % CA 总线输出
  \draw[hw bus] (10.7,1.0) -- (10.7,0.4);
  \node[hw label, font=\tiny, anchor=north] at (10.7,0.32) {CA 总线 $\to$ DRAM};
\end{pdfdiagram}
\diagramnote{bank 状态表与约束计数器。每个 bank 一行：状态、打开的行号、按 tRCD/tRAS/tRP 排列的一列倒计数器；发命令时装入对应时序参数，此后每拍减一，计数器为 0 是约束解除。就绪位针对队列中指向该 bank 的请求逐一判定（图中 bank 0/1 对应读请求，bank 2 对应激活请求），判定是纯组合逻辑，几十个请求的 ready 位同拍产生，优先级仲裁一拍内选出下一条命令，编码器把它翻译成 CA 总线上的控制线组合。bank 3 的 tRP 计数器未归零，本拍只能等待。}

这张图把"控制"二字展开到了电路尺度：内存控制器不知道也不关心什么是纳秒，它只数自己的时钟。DRAM 数据手册里的纳秒参数，在上电时被固件换算成拍数写进控制器寄存器——这本身就是一次软件对硬件的"参数配置"，与 1.4 节给 8254 装初值没有本质区别。

\topic{2.3\quad 内存训练：控制器如何做系统辨识}

上电瞬间，内存控制器面对的是一个参数未知的对象。DRAM 的时序参数写在 SPD（内存条上的参数 EEPROM，见 3.6 与 4.5 节）里可以读出来，但信号从控制器引脚到 DRAM 引脚要走多快——这取决于这块主板的走线长度、这根 DIMM（Dual In-line Memory Module，双列直插内存条）的插槽位置、今天的温度——没有任何地方写着。存储器件结构章第三节已经逐项讲过训练的操作细节：写均衡扫 DQS-CK 相位、读 per-bit 去偏斜为每根 DQ 扫描延迟线找出 PASS 窗口取中点、VREF 训练把时间与电压合成二维眼图。本节把这些步骤抽象回控制问题，因为它们是本章反馈主题的第一个完整实例。

训练的第一步是\hwkey{测量}，而测量的信息通道窄得惊人：写均衡时 DRAM（x8 器件）只能通过 DQ[0] 一根线、每拍回报 1 个 bit——DQS 边沿在 CK 之前还是之后。控制器就用这 1 bit 做二分或步进搜索，逐拍调整 DQS 输出延迟，直到回报翻转。第二步是\hwkey{参数搜索}：读去偏斜对每根 DQ 独立扫描延迟线的全部档位，记录 PASS 区间的左右边界，取中心——搜索目标是让采样点离两个失效边界等距，即最大化噪声裕量。第三步是\hwkey{写回配置}：搜索到的几十个最优值（每根 DQ 的延迟、每个 VREF 档位）写进 PHY 的配置寄存器。注意这正是 1.1 节的模式：PHY 是一台状态机，训练结果就是它的控制寄存器值，只是这次写寄存器的"软件"是固件里的内存参考代码（MRC），而不是操作系统。

把三步连起来看：发激励（写已知 pattern）、测响应（回读比较）、按误差调整参数、直到误差最小——这与磁盘伺服用 PES 调整音圈电流（存储器件结构章第六节）、与锁相环用相位差调整振荡频率，是同一个"测量—比较—调整"结构。差别只在环的速度：训练是固件闭合的慢环，一次冷启动执行一次，耗时几十到几百毫秒；训练后 PHY 按这组静态参数工作，直到下一次重训。而且慢环并不止于启动——LPDDR 平台（及部分 DDR5 平台）支持运行期周期性重训，跟踪温度与电压漂移；JEDEC DDR5 规定的是周期性 ZQ 校准，因为被控对象（走线延迟）本身随环境缓慢变化。本章第五节会把"测量—比较—调整"形式化成反馈控制的一般结构，内存训练、磁盘伺服、时钟恢复都将成为它的特例。
% ch06b 第三节：平台控制器
\section{三、平台控制器：从芯片组到 IOMMU}

经典芯片与平台章已经把芯片组迁移的轮廓画过一遍：分立的 8237/8259/8254 先被收进北桥与南桥，内存控制器随后搬进 CPU，南桥余部改称 PCH，最终走向 SoC。那一章回答的是“搬了什么”；本节回答“搬完之后它们怎么工作”：设备如何被发现并领到地址窗口，中断如何从一根引脚演化成一次内存写，DMA 地址如何被再翻译一次，操作系统的心跳由哪颗定时器发出，以及主板上那些不起眼却不能缺席的小控制器各管什么。

\begin{keyidea}
平台控制器的演进有一条统一的主线：\hwkey{能用地址空间表达的，就不用专用引脚}。设备配置空间被映射成一段 MMIO（ECAM），读设备标识就是一次普通 load；中断从 IRQ 引脚变成设备往约定地址的一次写（MSI）；DMA 地址治理借用了与 MMU 相同的页表结构（IOMMU）。整机的控制平面因此从“一堆专用信号线”统一成“地址空间中的读写事务”——这正是主存与总线章里 MMIO 思想在平台尺度上的展开。
\end{keyidea}

%======================================================================
\topic{3.1\quad 演进：北桥/南桥、PCH 与 SoC 各管什么}
%======================================================================

南北桥时代的分工按“离 CPU 多近、要多快”切分。\hwkey{北桥}（northbridge）坐在 CPU 前端总线上，管三条最快的路径：内存控制器（DRAM 命令与调度）、图形接口（AGP，Accelerated Graphics Port，加速图形端口；后来是 PCIe $\times$16）和 cache/前端总线一致性。\hwkey{南桥}（southbridge）管全部低速 I/O：IDE/SATA、USB、PCI、LPC（Low Pin Count，低引脚数总线），并把 8259 中断控制器、8237 DMA、8254 定时器、RTC 这些经典芯片的功能各收一个副本进片内，外加电源管理与 ACPI。两桥之间用一条专用 hub 链路相连，南桥的全部流量都要经过北桥才能到达内存。

把内存控制器从北桥搬进 CPU 是收益最大的一次迁移，原因要分两层看。第一层是延迟：板级走线上信号每走 10 cm 约需 \hwtiming{0.67 ns} 飞行时间（电源、时钟与信号完整性章第三节给出的量级），再算上两端的收发缓冲与仲裁，一次片外往返在纳秒到十几纳秒量级；北桥时代每次 cache miss 都要先走“CPU$\to$北桥”这一程再去 DRAM，比控制器片内方案多出几十纳秒。取指和读数是整机最频繁的路径，这段开销被乘以全部访存次数。第二层是功耗：驱动板级走线要把 50 $\Omega$ 量级的负载在每比特内反复充放，片外 I/O 的能耗典型在 \hwtiming{10--30 pJ/bit}，片内互连则低于 1 pJ/bit——同一份数据少出一次芯片，就省下一次数量级的能量。

迁移分三步走完，每一步搬走的东西都很具体：

\begin{longtable}{L{0.14\linewidth}L{0.40\linewidth}L{0.36\linewidth}}
\toprule
阶段 & 结构 & 被搬近的路径 \\
\midrule
北桥 + 南桥 & 北桥管内存/图形，南桥管低速 I/O，hub 链路相连 & 内存与图形集中到北桥，I/O 集中到南桥 \\
PCH & 内存控制器、PCIe 根复合体进 CPU；南桥余部 + 固件接口改称 PCH（Platform Controller Hub），经 DMI（Direct Media Interface，直接媒体接口）连 CPU & cache miss 到 DRAM 不再跨芯片；DMI 本质是一条窄化的 PCIe（$\times$4，量级 4--8 GB/s） \\
SoC & 图形、显示、外设控制全部收进同一 die 或同一封装 & 常见路径不出芯片，只剩电源、时钟和低速引脚外连 \\
\bottomrule
\end{longtable}

既然集成度提高总是更快更省，为什么台式机至今仍是“CPU + PCH”两颗？三个原因都很实际。其一是引脚：低速 I/O 数量大、各占一根引脚，收进 CPU 要占用大量引脚与 die 边缘面积，换来的收益却很小——SATA 盘本来就比 DRAM 慢好几个数量级。其二是工艺：LPC、RTC、模拟监控这类电路工作电压高、对最先进逻辑工艺并不友好，放在成熟工艺的 PCH 上更便宜。其三是产品分档：同一颗 CPU 搭配不同档次的 PCH（通道数、USB 口数不同），不必为每个市场单独流片。理解“什么该集成”与理解“集成了什么”同样重要：集成的是延迟敏感的高频路径，留下的是引脚密集、速度不敏感的长尾。

\begin{pdfdiagram}[x=1cm,y=1cm]
  % === CPU die 容器 ===
  \draw[draw=BookNavy, fill=BookCodeBg, line width=0.6pt, rounded corners=4pt] (0.5,8.2) rectangle (6.4,10.6);
  \node[hw label, font=\scriptsize\bfseries, text=BookNavy, anchor=west] at (0.7,10.3) {CPU die};
  \node[hw compute] (core) at (1.5,9.4) {核心+LLC};
  \node[hw control] (rc) at (3.45,9.4) {PCIe 根\\复合体};
  \node[hw compute] (imc) at (5.35,9.4) {内存控制器};
  % === DIMM ===
  \node[hw memory] (dimm) at (8.3,9.4) {DDR5 DIMM\\$\times$2};
  \draw[hw bus] (imc.east) -- (dimm.west);
  \node[hw label, font=\tiny, anchor=south] at (8.3,10.05) {DDR5-4800 双通道\\$\approx$77 GB/s};
  % === PCIe 设备 ===
  \node[hw compute] (gpu) at (1.5,6.2) {GPU\\$\times$16};
  \node[hw compute] (nvme) at (4.0,6.2) {NVMe SSD\\$\times$4};
  \draw[BookMuted, line width=0.62pt] (rc.south) -- (3.45,7.7);
  \draw[BookMuted, line width=0.62pt] (1.5,7.7) -- (4.0,7.7);
  \draw[hw bus] (1.5,7.7) -- (gpu.north);
  \draw[hw bus] (4.0,7.7) -- (nvme.north);
  \node[hw label, font=\tiny, anchor=west] at (1.65,7.15) {$\approx$32 GB/s};
  \node[hw label, font=\tiny, anchor=west] at (4.15,7.15) {$\approx$8 GB/s};
  % === DMI 与 PCH ===
  \draw[hw bus] (5.9,8.2) -- (5.9,5.0);
  \node[hw label, font=\tiny, anchor=west] at (6.05,6.5) {DMI $\times$4 $\approx$8 GB/s};
  \draw[draw=BookNavy, fill=white, line width=0.6pt, rounded corners=4pt] (1.0,2.8) rectangle (6.8,5.0);
  \node[hw label, font=\scriptsize\bfseries, text=BookNavy, anchor=west] at (1.2,4.68) {PCH};
  \node[hw label, text=BookInk] at (3.9,3.9) {SATA · USB 主控 · 以太网 MAC\\[2pt] LPC/eSPI · SPI · SMBus\\[2pt] RTC · GPIO · 定时器 · 中断};
  % === PCH 下游设备 ===
  \draw[BookMuted, line width=0.62pt] (3.9,2.8) -- (3.9,2.2);
  \draw[BookMuted, line width=0.62pt] (1.05,2.2) -- (8.4,2.2);
  \node[hw memory] (sata) at (1.05,1.1) {SATA 盘};
  \node[hw state] (usb) at (2.95,1.1) {USB 设备};
  \node[hw state] (gbe) at (4.65,1.1) {GbE PHY};
  \node[hw compute] (ec) at (6.45,1.1) {EC /\\Super I/O};
  \node[hw memory] (flash) at (8.4,1.1) {BIOS\\flash};
  \draw[hw bus] (1.05,2.2) -- (sata.north);
  \draw[hw bus] (2.95,2.2) -- (usb.north);
  \draw[hw bus] (4.65,2.2) -- (gbe.north);
  \draw[hw bus] (6.45,2.2) -- (ec.north);
  \draw[hw bus] (8.4,2.2) -- (flash.north);
  \node[hw label, font=\tiny, anchor=west, fill=white, inner sep=1pt] at (1.15,1.88) {6 Gb/s};
  \node[hw label, font=\tiny, anchor=west, fill=white, inner sep=1pt] at (3.05,1.88) {10/20 Gb/s};
  \node[hw label, font=\tiny, anchor=west, fill=white, inner sep=1pt] at (4.75,1.88) {1 Gb/s};
  \node[hw label, font=\tiny, anchor=west, fill=white, inner sep=1pt] at (6.55,1.88) {LPC/eSPI};
  \node[hw label, font=\tiny, anchor=west, fill=white, inner sep=1pt] at (8.5,1.88) {SPI $\approx$50 MB/s};
\end{pdfdiagram}
\diagramnote{现代 PC 平台拓扑与带宽量级（单向）。CPU die 内的内存控制器与 PCIe 根复合体接管最快的路径：DRAM 直连内存控制器，GPU 与 NVMe SSD 挂在根复合体上；PCH 经 DMI 汇聚全部低速 I/O。注意图上每一处带宽差距都在 10 倍以上——DRAM 约 77 GB/s，DMI 约 8 GB/s，SATA 约 600 MB/s，LPC 与 SPI 只有数十 MB/s——这个落差正是“快慢分桥”的原因。软件侧的直接对应：\code{lspci -tv} 输出的顶层设备 host bridge 与 ISA bridge，就是 CPU 内根复合体与 PCH 在 PCI 模型中的投影。}

%======================================================================
\topic{3.2\quad PCIe 根复合体与设备枚举}
%======================================================================

\term{根复合体}{root complex}是 CPU 与 PCIe 之间的翻译器，工作方向有两个。向下，它把 CPU 对物理地址的读写翻译成 PCIe 事务层包（TLP）发往相应设备——对 BAR 窗口的 MMIO 写变成一个 Memory Write TLP；向上，它把设备发来的 TLP（DMA 读写、MSI 中断消息）翻译成对 CPU 一致性域内存的访问。路由规则分两类：\hwkey{地址路由}用于内存事务，按目标地址落在哪个桥或设备的窗口里逐级转发；\hwkey{ID 路由}用于配置事务与消息，按 BDF 定位。所谓 \term{BDF}{bus:device:function} 是 PCIe 给每个功能的编号：8 位总线号、5 位设备号、3 位功能号，合起来就是 \code{lspci} 输出里 \code{0000:03:00.0} 这样的地址。

设备要被软件使用，先要能被软件找到。每个 PCIe 功能带 4 KiB \hwkey{配置空间}，头部 64 字节是标准布局：偏移 \code{0x00} 处是厂商 ID 与设备 ID，读回全 1（\code{0xFFFFFFFF}）表示该位置没有设备；偏移 \code{0x10} 起是 BAR。配置空间本身又通过 \term{ECAM}{enhanced configuration access mechanism} 映射进物理地址空间：把 256 条总线 $\times$ 32 个设备 $\times$ 8 个功能 $\times$ 4 KiB 共 256 MiB 映射到一段 MMIO，某个寄存器的地址为

\[
\text{addr} = \text{base} + (\text{bus}\ll 20) \mathbin{|} (\text{dev}\ll 15) \mathbin{|} (\text{fun}\ll 12) + \text{reg}
\]

基址 base 由固件经 ACPI 的 MCFG 表告诉操作系统。于是“读一块网卡的设备 ID”对内核来说就是一次普通 load——配置这件事不需要任何专用指令，这是本节核心思想的第一处体现。

配置头分两种。类型 0 是端点设备（网卡、SSD、显卡），类型 1 是桥（根端口、switch 下游口），区别在于中段：端点用 BAR 声明自己要多少地址窗口，桥用三个总线号寄存器声明自己下游挂着哪段总线。

\begin{pdfdiagram}[x=1cm,y=1cm]
  \node[hw label, font=\scriptsize\bfseries, text=BookNavy] at (3.0,5.35) {类型 0：端点设备};
  \node[hw label, font=\scriptsize\bfseries, text=BookNavy] at (9.0,5.35) {类型 1：桥};
  % 左列：类型 0（0x0C 行单独画，便于高亮）
  \foreach \y/\off/\txt in {
    4.60/0x00/Device ID（16 位）｜Vendor ID（16 位）,
    3.88/0x04/Status｜Command,
    2.44/0x10/BAR0 … BAR5（0x10--0x24，六个窗口）,
    1.72/0x2C/子系统 Vendor ID｜子系统 Device ID,
    1.00/0x3C/Int Pin｜Int Line} {
    \draw[fill=BookCodeBg, draw=BookRule, line width=0.4pt] (0.6,\y-0.3) rectangle (5.4,\y+0.3);
    \node[hw label, font=\tiny, text=BookInk] at (3.0,\y) {\txt};
    \node[hw label, font=\tiny, anchor=east] at (0.55,\y) {\off};
  }
  \draw[fill=BookAmber!20, draw=BookRule, line width=0.4pt] (0.6,2.86) rectangle (5.4,3.46);
  \node[hw label, font=\tiny, text=BookInk] at (3.0,3.16) {Header Type = 0x00（端点）};
  \node[hw label, font=\tiny, anchor=east] at (0.55,3.16) {0x0C};
  % 右列：类型 1
  \foreach \y/\off/\txt in {
    4.60/0x00/Device ID｜Vendor ID,
    3.88/0x04/Status｜Command,
    2.44/0x18/Primary｜Secondary｜Subordinate 总线号,
    1.72/0x20/下游内存窗口 基址/限长} {
    \draw[fill=BookCodeBg, draw=BookRule, line width=0.4pt] (6.6,\y-0.3) rectangle (11.4,\y+0.3);
    \node[hw label, font=\tiny, text=BookInk] at (9.0,\y) {\txt};
    \node[hw label, font=\tiny, anchor=east] at (6.55,\y) {\off};
  }
  \draw[fill=BookAmber!20, draw=BookRule, line width=0.4pt] (6.6,2.86) rectangle (11.4,3.46);
  \node[hw label, font=\tiny, text=BookInk] at (9.0,3.16) {Header Type = 0x01（桥）};
  \node[hw label, font=\tiny, anchor=east] at (6.55,3.16) {0x0C};
\end{pdfdiagram}
\diagramnote{PCIe 配置头对比（每格一行为一个 32 位寄存器，左侧标注偏移）。前 16 字节两种类型相同：Vendor/Device ID 标识身份，Command 寄存器里的 Memory Space Enable 与 Bus Master Enable 两位决定设备是否响应 MMIO、是否允许发起 DMA。高亮行是 Header Type 字段（位于字节 0x0E，图中按所在行首偏移 0x0C 标注），枚举软件靠它分流。类型 0 的六个 BAR 声明地址窗口需求；类型 1 的三个总线号寄存器圈出下游总线范围，上游按它做 ID 路由。}

BAR（Base Address Register，基址寄存器）的分配过程是一段精巧的硬件约定。上电时 BAR 内容未定，固件按三步处理每个 BAR：先写入全 1 再读回——窗口位可写、低位硬连线为 0，读回值的最低置位即窗口大小，例如读回 \code{0xFFFF0000} 表示设备要 64 KiB；然后从空闲物理地址池中分一段满足对齐的窗口写回 BAR；最后置 Command 寄存器的使能位。此后设备内部的地址译码器只响应落在该窗口内的访问。64 位 BAR 占相邻两个槽位；最低位区分内存窗口与 I/O 端口窗口。NVMe 控制器、网卡、显卡的驱动能看到的全部 MMIO 寄存器，都住在这样分出来的窗口里。

枚举（enumeration）把上述过程组织成一次深度优先遍历。固件从根总线 0 开始扫描全部 dev/fun 位置，读到桥就给它分配一个 secondary 总线号并递归扫描下游，回溯时把下游最大总线号写进 subordinate 寄存器；读到端点就探测 BAR 并分配窗口：

\begin{lstlisting}
scan_bus(bus):
    for dev in 0..31, fun in 0..7:
        id = cfg_read(bus, dev, fun, 0x00)
        if id == 0xFFFFFFFF: continue     # 该位置无设备
        if header_type == 0x01:           # 桥
            set secondary = next_bus++
            scan_bus(secondary)           # 递归下游
            set subordinate               # 回填下游最大总线号
        else:                             # 端点
            probe BARs                    # 写全 1 探测大小
            分配 MMIO/IO 窗口并写回 BAR
\end{lstlisting}

遍历结束时，整棵树的每个功能都有了确定的 BDF、确定的地址窗口，桥上的总线号范围支持 ID 路由，整条链路可以开始跑内存事务。这段工作在传统上由 BIOS/UEFI 固件完成，操作系统启动后通常沿用固件分配的结果（Linux 也保留重新分配的能力）；\code{/sys/bus/pci/devices/} 下的目录名就是 BDF，\code{resource} 文件里就是每个 BAR 领到的窗口。

\begin{warningbox}
BAR 窗口是“一段地址”，不是“一段内存”。这段物理地址背后没有 DRAM，读它得到的是设备寄存器的当前值，写它驱动的是设备的状态机——主存与总线章第三节讲的 MMIO 副作用规则在这里全部适用：窗口必须按不可缓存（UC）属性映射，不能用 \code{memcpy} 式的块拷贝访问 32 位寄存器，也不能指望写返回就代表设备动作完成。另外，窗口之间、窗口与 DRAM 之间不允许重叠；重叠时两个目标同时响应一次访问，结果是总线冲突或数据错乱，且往往表现为与软件无关的随机故障。
\end{warningbox}

%======================================================================
\topic{3.3\quad 中断控制器演进：8259、APIC 与 MSI/MSI-X}
%======================================================================

8259 的内部机制——IRR/IMR/ISR 三张寄存器、EOI（End Of Interrupt，中断结束）纪律、优先级裁决——主存与总线章第四节已经讲过，经典芯片与平台章也走完了一次 INTR 从请求到 INTA 周期的完整链路。这里只补它在整机中的位置与瓶颈。PC/AT 起用两片 8259 级联：从片的 INT 输出接主片的 IR2，共提供 15 条可用中断线，IRQ 号与设备的对应关系被写死成平台约定：

\begin{longtable}{L{0.10\linewidth}L{0.34\linewidth}L{0.49\linewidth}}
\toprule
IRQ & 固定对应 & 说明 \\
\midrule
0 & PIT 8254 通道 0 & 系统心跳，见 3.5 小节 \\
1 & 键盘控制器 & 8042，后由 EC 继承 \\
2 & 级联线 & 从片 INT 接主片 IR2，本身不可用 \\
3 / 4 & COM2 / COM1 & 串口 \\
6 & 软盘控制器 & 与 DMA 通道 2 配套 \\
8 & RTC & 闹钟与周期中断 \\
12 & PS/2 鼠标 & \\
14 / 15 & 主/从 IDE 通道 & \\
\bottomrule
\end{longtable}

这套结构在单 CPU、十几个设备的时代够用，瓶颈在四个方向同时出现：线数封顶 15，设备多起来只能共享；共享线在边沿触发下不可靠，两个设备恰好同时请求时后一个沿会被吞掉；优先级固定，IRQ 号小的永远压过大的；而且整套机制只认一个 CPU——向量表、EOI、屏蔽全是全局唯一的一份。多核时代需要的恰恰相反：中断线要多、要能定向到指定核、要让高频设备（网卡、SSD）独占向量。

APIC（Advanced Programmable Interrupt Controller，高级可编程中断控制器）把中断控制器拆成两半。\hwkey{Local APIC} 每个 CPU 核一个，MMIO 基址 \code{0xFEE00000}，内部有任务优先级寄存器 TPR（只有高于该阈值的中断才被投递）、按 256 个向量各一位的 IRR/ISR、EOI 寄存器、本地向量表 LVT（本核的定时器、温度传感器、性能计数溢出等本地中断源）以及标识自己身份的 APIC ID。\hwkey{IO APIC} 在 PCH 内，对外收设备的 IRQ 引脚，内部是一张重定向表（典型 24 项），每项指定：这个引脚来的请求译成几号向量、按什么方式投递（固定、最低优先级、SMI 系统管理中断、NMI）、送往哪个或哪些核的 APIC ID、电平还是边沿触发、极性与屏蔽。设备拉起 IRQ 线后，IO APIC 查表生成一条中断消息，经系统互连送到目标核的 Local APIC，后者按 TPR 与 IRR/ISR 判优，在指令边界向 CPU 投递向量号，CPU 查 IDT 进入处理程序。\hwkey{中断亲和性}（affinity）就是写重定向表的目标字段：把网卡队列的中断固定在某个核上，cache 与协议栈状态都留在本地；Linux 的 irqbalance 守护进程做的就是动态调整这些字段。

\term{MSI}{message-signaled interrupt} 更进一步，干脆取消中断引脚：设备要发中断时，往 \code{0xFEExxxxx} 区间做一次 32 位内存写——地址位编码目标核的 APIC ID，数据低 8 位是向量号。这次写经 PCIe 到达根复合体，被解释成一条中断消息转交 Local APIC。预先约定的那对“地址 + 数据”由操作系统在设备初始化时写进设备的 MSI capability 寄存器。MSI 允许每设备最多 32 个向量且必须取 2 的幂连续分配；\term{MSI-X}{MSI extended} 把这个限制放宽到 2048 个向量，每向量有独立的地址、数据与屏蔽位，表格本身放在 BAR 窗口里。MSI-X 的意义不只是数量：网卡、NVMe SSD 可以给每个收发队列配一个独立向量、定向到不同核，中断路径与数据路径按队列一一对应——存储器件结构章第五节讲的 NVMe 每核一对队列，配对的正是一个 MSI-X 向量。

MSI 还有一条容易被忽视的保序性质：它是一次 posted write，遵循 PCIe 的排序规则，不会越过同一设备先发的 DMA 数据写。因此“DMA 写完数据 $\to$ 发 MSI”这个序列到达内存系统的顺序与发出顺序一致，中断处理程序读到的一定是完整数据，不需要额外的人工同步——引脚时代“中断先到、数据还在路上”的竞态在机制上被消除了。

\begin{pdfdiagram}[x=1cm,y=1cm]
  \node[hw compute] (legacy) at (1.3,5.8) {传统设备\\（IRQ 引脚）};
  \node[hw compute] (pciedev) at (1.3,2.8) {PCIe 设备\\（MSI/MSI-X）};
  \node[hw control] (ioapic) at (5.0,5.8) {IO APIC（PCH 内）\\24 项重定向表};
  \node[hw control] (lapic) at (9.9,5.8) {Local APIC（每核一个）\\\texttt{0xFEE0\,0000}};
  \node[hw state] (cpu) at (9.9,7.9) {CPU 核\\IDT $\to$ 处理程序};
  \draw[hw bus] (legacy.east) -- (ioapic.west);
  \node[hw label, font=\tiny, anchor=south] at (2.8,5.92) {IRQn};
  \draw[hw bus] (ioapic.east) -- (lapic.west);
  \node[hw label, font=\tiny, anchor=south] at (7.45,5.92) {中断消息};
  \draw[hw bus] (pciedev.east) -- (9.9,2.8) -- (lapic.south);
  \node[hw label, font=\tiny, anchor=south, fill=white, inner sep=1pt] at (5.9,3.02) {MSI：写 \texttt{0xFEEx\,xxxx}，地址选核、数据含向量};
  \draw[hw bus] (lapic.north) -- (cpu.south);
  \node[hw label, font=\tiny, anchor=east] at (9.7,6.85) {按 TPR 判优后投递};
\end{pdfdiagram}
\diagramnote{中断从设备到 CPU 的两条投递路径在 Local APIC 汇合。引脚路径：设备拉 IRQ 线，IO APIC 查重定向表得到向量号与目标核，生成中断消息送往该核的 Local APIC。MSI 路径：设备直接对 \code{0xFEE} 开头的地址做一次内存写，地址本身携带目标核，数据携带向量号，不经任何引脚。Local APIC 按 TPR 与 IRR/ISR 判优后在指令边界投递，CPU 查 IDT 进入处理程序。软件侧的观察窗口是 \code{/proc/interrupts}：每一行是一个向量，每一列是一个核的计数——从计数的分布可以直接读出亲和性配置与设备申请到的 MSI-X 向量数。}

\begin{classicnote}
CSAPP 第 8 章把中断归入“异步异常”，并用定时器与磁盘完成通知作为例子说明控制流如何在任意指令边界被打断。本节给出的是“异步”两个字背后的投递硬件：从 8259 的一根 INTR 引脚，到 IO APIC 的一张重定向表，再到 MSI 的一次内存写——异步性没有变，变的是路由的表达能力。
\end{classicnote}

%======================================================================
\topic{3.4\quad DMA 控制器与 IOMMU}
%======================================================================

8237 作为 DMA 控制器的经典样本，主存与总线章第四节已经讲到通道寄存器、页寄存器与单字/块模式的细节，这里只补它的整机编年：PC/AT 把两片 8237 级联出 7 个通道（通道 4 用于级联本身），软盘固定占通道 2，声卡驱动配置里的“DMA 1、DMA 5”说的就是它；它本质是一台独立的搬运状态机，每传一字节都要向 CPU 借一次总线（单字模式）。1990 年代起它被淘汰的原因很直接——速度。每字节一次握手让它的吞吐停在 MB/s 量级，而 PCI 设备自己发起突发传输能快两个数量级。于是 DMA 引擎改为分散在每个高速设备内部，即 \hwkey{bus-master DMA}：设备自己当总线主，按描述符环批量搬运（描述符环与可见性规则见主存与总线章，NVMe 的队列与 PRP 见存储器件结构章第五节）。“中央 DMA 控制器”这个角色就此消失，8237 只在 LPC 上留一个兼容影子。

但 DMA 引擎分散之后，留下一个集中式时代不突出的问题：\hwkey{地址治理}。设备发起的 DMA 读写不经过 CPU，也不经过 MMU——页表管的是指令发出的虚拟地址，管不到设备发出的地址。一块固件有缺陷的网卡、一个有 bug 的驱动，往描述符里写了错误地址，DMA 就原样写坏任意物理内存，连页错误都不会产生。虚拟化场景更尖锐：把网卡直通给虚拟机后，客户机驱动按它以为的物理地址（guest 物理地址）编程设备，而这些地址在主机上对应着别的内存。

\term{IOMMU}{I/O memory management unit}（Intel 称 VT-d，AMD 称 AMD-Vi）的解法是把 MMU 的页表机制复制一份给设备用。每个 DMA 事务都携带来源标识 source-id，就是发起者的 BDF；IOMMU 先用总线号索引\hwkey{根表}（256 项），再用设备号/功能号索引\hwkey{上下文表}，得到该设备所属域的二级页表基址；随后像 MMU 一样做多级页表 walk，把设备发出的 \term{IOVA}{I/O virtual address} 翻译成真正的物理地址，翻译结果缓存于 IOTLB，高端设备还可以持有 ATS（Address Translation Services，地址翻译服务）缓存把翻译移到设备侧。查不到映射或权限不符，本次 DMA 被拒绝并产生 IOMMU fault 中断，错误被限制在事件级别而不是内存损坏。

\begin{pdfdiagram}[x=1cm,y=1cm]
  \node[hw compute] (dev) at (1.0,3.0) {PCIe 设备\\（bus-master）};
  \node[hw state] (root) at (3.6,3.0) {根表\\按 bus 索引};
  \node[hw state] (ctx) at (6.2,3.0) {上下文表\\按 dev/fun};
  \node[hw state] (spt) at (8.6,3.0) {二级页表\\IOVA $\to$ PA};
  \node[hw memory] (dram) at (10.7,3.0) {DRAM\\物理页};
  \draw[hw bus] (dev.east) -- (root.west);
  \draw[hw bus] (root.east) -- (ctx.west);
  \draw[hw bus] (ctx.east) -- (spt.west);
  \draw[hw bus] (spt.east) -- (dram.west);
  \node[BookRed, font=\tiny\bfseries] at (2.3,3.3) {\textcircled{1}};
  \node[BookRed, font=\tiny\bfseries] at (4.9,3.3) {\textcircled{2}};
  \node[BookRed, font=\tiny\bfseries] at (7.4,3.3) {\textcircled{3}};
  \node[BookRed, font=\tiny\bfseries] at (9.8,3.3) {\textcircled{4}};
  \node[hw label, font=\tiny, text=BookRed] at (6.0,2.1) {任何一级查无映射或权限不符 $\to$ IOMMU fault，本次 DMA 被拒绝};
  \node[hw label, font=\tiny] at (6.0,1.72) {OS 经 DMA 映射接口为每个缓冲区建立 IOVA $\to$ PA 页表项，根表基址由 OS 写入};
\end{pdfdiagram}
\diagramnote{IOMMU 对一次 DMA 的翻译链。\textcircled{1} 设备发起 DMA，事务携带来源标识（BDF）与 IOVA；\textcircled{2} 根表项（按总线号索引，共 256 项）给出上下文表基址；\textcircled{3} 上下文项（按设备号/功能号索引）给出该设备域的二级页表基址与域标识；\textcircled{4} 页表 walk（结果缓存于 IOTLB）得出物理地址，访问发往 DRAM。结构上它与 MMU 完全同构——区别只在触发者是设备而不是指令，页表按“设备域”而不是按“进程”组织。}

IOMMU 的用途按重要性排有三层。第一是保护：操作系统只为缓冲区建立映射，设备越出映射范围即 fault，驱动 bug 与恶意固件的破坏半径被限制在分配给它的页内。第二是虚拟化直通：把物理网卡分给虚拟机，客户机驱动照常工作，IOMMU 把 guest 物理地址再翻译一层到 host 物理地址——虚拟机获得接近原生的 I/O 性能，隔离性却由硬件维持，云主机的网卡直通走的就是这条路。第三是兼容：只能发 32 位地址的老设备，经重映射后可以使用任意位置的缓冲区，不再依赖低 4 GiB 的保留区。代价同样实在：映射的建立与撤销、IOTLB 失效广播都有开销，高频小包场景下 DMA 映射开销可观，工程上的对策是大页、预分配映射池，以及可信环境下用直通模式（Linux 的 \code{iommu=pt}）让 IOVA 恒等映射物理地址、跳过逐缓冲区建表。

\begin{warningbox}
开启 IOMMU 之后，“DMA 地址等于物理地址”这条假设就失效了。驱动拿到的 \code{dma_addr_t} 是 IOVA，只有通过 DMA 映射接口（\code{dma_map_single} 等）建立映射后得到的地址才能写给设备；把 \code{virt_to_phys} 的结果直接填进描述符，轻则在 IOMMU 开启的机器上触发 fault，重则在映射恰好撞上别的页时写坏无关内存。这也是跨平台驱动代码必须走 DMA API 而不能省的原因。
\end{warningbox}

%======================================================================
\topic{3.5\quad 定时器家族：操作系统的心跳从哪来}
%======================================================================

操作系统对时间有三种需求：周期性的节拍（驱动调度、时间片、超时轮询）、到点即响的单次截止（高精度定时器）、以及可读取的时间戳（测量与日志）。一颗定时器很难同时做好三件事，于是平台上并存着一整族定时器，各管一段。

\hwkey{PIT 8254} 是最老的一员。输入时钟 \hwtiming{1.193182 MHz}（14.31818 MHz 晶振的十二分频，这个奇怪的数字来自 PC 对 NTSC 彩色副载波的复用），片内三个 16 位通道：通道 0 接 IRQ0 产生周期中断，通道 1 曾用于 DRAM 刷新，通道 2 驱动蜂鸣器。除数最大 65536，对应最低频率约 \hwtiming{18.2 Hz}——早期 DOS 的心跳就是它；后来的操作系统把节拍提到 100 Hz、250 Hz、1000 Hz。它的局限在多核时代很突出：全系统只有一颗，重编程要走慢速端口 I/O，分辨率也只有约 0.84 $\mu$s。

\hwkey{HPET} 是替代方案：一个 64 位主计数器以不低于 \hwtiming{10 MHz} 的频率走字（典型值 14.31818 MHz），配至少 3 个比较器，任一比较器匹配即可产生中断，支持单次与周期两种模式。寄存器全部 MMIO 映射，读写比端口 I/O 快，比较器数量也够多个定时源并存。但它仍在 PCH 里，读一次主计数器是一次跨芯片的 MMIO 往返，代价在数十纳秒量级——做中断源合格，做高频时间戳不合格。

\hwkey{Local APIC timer} 把定时器挪回每个核内部：以总线或核心时钟分频后的频率对一个初始计数值做减计数，减到 0 经 LVT 投递中断，有单次、周期、TSC-deadline 三种模式。每核私有意味着中断天然定向本核、不需要任何跨核仲裁，这正是无滴答（tickless）内核给每个核独立编程下一次唤醒点的硬件基础；要注意的是深度睡眠状态可能停走这颗定时器，唤醒源需要另行安排。

\hwkey{TSC}（Time Stamp Counter，时间戳计数器）严格说不是定时器而是计数器：每核一个 64 位寄存器，随核心基准频率走字，现代处理器保证 invariant——变频、休眠都不改变它的速率与连续性。\code{rdtsc} 读它一次约 \hwtiming{20--30 拍}，分辨率在纳秒量级，因此成为时间戳的事实标准：\code{clock_gettime} 的快路径（vDSO，virtual dynamic shared object，内核映射到用户态的共享页）就是读 TSC 再乘一个校准好的换算系数，根本不进内核。TSC-deadline 模式则把 Local APIC timer 与 TSC 接在一起：直接写一个目标 TSC 值作为比较点，省去了“想要 5 $\mu$s 后中断”到“该写多少计数”之间的频率换算。

\begin{longtable}{L{0.13\linewidth}L{0.19\linewidth}L{0.18\linewidth}L{0.17\linewidth}L{0.21\linewidth}}
\toprule
器件 & 时钟源 & 分辨率/读代价 & 中断去向 & 典型软件用途 \\
\midrule
PIT 8254 & 1.193182 MHz & $\sim$0.84 $\mu$s/端口 I/O & 全局 IRQ0 & 早期 tick，今多为兼容保留 \\
HPET & $\geq$10 MHz（典型 14.3 MHz） & $\sim$70 ns/跨芯片 MMIO & 可经 IO APIC 定向 & 传统高精度 tick、校准参考 \\
LAPIC timer & 总线/核心时钟分频 & ns 级/片内寄存器 & 本核 LVT & tickless 内核的 per-cpu 定时 \\
TSC & invariant 核心基准频率 & ns 级/$\sim$20--30 拍 & 无（仅计数） & 时间戳、vDSO 快路径、性能测量 \\
RTC（见 3.6） & 32.768 kHz 音叉晶振 & 1 s & IRQ8（闹钟） & 墙钟初值、定时开机 \\
\bottomrule
\end{longtable}

把这张表与“操作系统的心跳”对应起来：早期内核的 tick 来自 PIT 或 HPET 的周期中断，每个 tick 到来时调度器统计时间消耗、检查时间片、推进超时队列；现代 tickless 内核把周期中断省掉，空闲核按最近的截止时刻给 LAPIC timer 编程一个单次比较，真正有事才醒。值得分清的是两组“精度”：PIT、HPET、LAPIC timer、TSC 四者的分辨率与访问代价相差三个数量级（PIT 的亚微秒到 TSC 的纳秒），但长期漂移都由晶振决定，全部在 ppm 量级（电源时钟章第二节讨论过晶振与抖动）——也就是说，这些定时器“走得一样准，看得不一样细”，跨天级别的准确度从来不是靠它们中的任何一个，而是靠 NTP 之类的协议持续校准。

%======================================================================
\topic{3.6\quad 小但不可缺的控制器：RTC、GPIO、SMBus、Super I/O、EC 与 TPM}
%======================================================================

平台上最后一类控制器速度慢、体积小，却各自守着一项整机缺它不可的职能。逐个看一遍，重点是它们各自在软件世界里留下的入口。

\hwkey{RTC}（Real-Time Clock，实时时钟）由一颗 32.768 kHz 音叉晶振驱动——选这个频率是因为 $32768=2^{15}$，十五级二分频恰好得到 1 Hz。它挂在纽扣电池上，整机关机后照走，寄存器里以 BCD 格式保存年月日时分秒，还能在设定时刻或按周期发闹钟中断（传统 IRQ8），旁带一小片同样由电池保持的 CMOS RAM，早年用来存启动顺序等固件设置。软件碰到它的方式很固定：内核启动早期读 RTC 初始化墙钟，之后系统时间靠 NTP 校准，RTC 只负责“下一次上电时还记得大概是几点”。它的精度在 $\pm$20 ppm 量级，一天能漂出秒级，所以角色是上电初值而不是计时基准——定时开机、从 S3/S5 唤醒闹钟才是它不可替代的部分。

\hwkey{GPIO}（General-Purpose I/O，通用输入输出）是 PCH 上一把可逐位配置的引脚：每位有方向、输出电平、输入状态三类寄存器，有的还带中断与去抖能力。电源键、复位线、LED、风扇转速检测、电平开关，整机上一切“非协议”的单比特信号都挂在 GPIO 上。软件侧，驱动直接读写 PCH 的 GPIO 寄存器块；Linux 的入口从 \code{/sys/class/gpio} 演进到了 libgpiod；慢速的器件协议（如某些传感器）也常靠软件翻转 GPIO 模拟出来。

\hwkey{I2C/SMBus} 的协议本身——两根开漏线、起始/停止条件、地址加读写位——主存与总线章第三节已经讲过；SMBus 是 I2C 在系统管理场景下的子集，速率 100 kHz，命令语义固定（读字节、写字节、块读）。它在整机中有两条著名用途。其一是内存条上的 \term{SPD}{Serial Presence Detect} EEPROM：固定占地址 \code{0x50}--\code{0x57}，里面存着厂商、容量、die 组织和一整套时序参数表，固件上电后第一件事就是沿 SMBus 逐条读出 SPD，据此配置内存控制器的频率与时序——存储器件结构章第二节那些 CL、tRCD、tRP 参数，开机时就是从这几百个字节里来的。其二是主板上的电压/温度传感器，lm-sensors 读到的数据全部来自这条总线。

\hwkey{Super I/O} 是 LPC 总线上的老牌合集芯片：串口 UART、并口、PS/2 键盘鼠标口、风扇控制与硬件监控各收一份，配置走 \code{0x2E}/\code{0x4E} 端口的索引/数据寄存器对。台式机主板上它至今常见，风扇调速曲线、电压监控都由它执行，内核的 hwmon 子系统里有专门的驱动家族。笔记本里这个角色通常并入了下一位。

\hwkey{EC}（Embedded Controller，嵌入式控制器）是笔记本主板上的一颗小单片机——历史上多为 8051 核，现在常见 ARM 或 RISC-V 核——接在常电域上，关机插电时也在运行。它管键盘矩阵扫描、电池充放电管理、风扇、温度采样、电源键与盖子开关，是“按一下电源键机器才会醒”这条链路的起点。与操作系统的接口走 ACPI：有事件时 EC 拉起 SCI（System Control Interrupt，系统控制中断）中断，OS 执行对应的 ACPI 控制方法，再经共享寄存器区读写 EC 状态。理解 EC 就理解了“关机不等于断电”：主板上始终有一台小电脑在值班，等待下一次按键。

\hwkey{TPM}（Trusted Platform Module，可信平台模块）是一颗独立的安全芯片（LPC、SPI 或 I2C 接口；现代 CPU 也提供固件实现 fTPM），内部有非对称与对称加密引擎、单调计数器、受保护的密钥存储，以及一组 \term{PCR}{platform configuration register}。PCR 的特性是只能 extend 不能写：新值 $=$ hash（旧值 $\|$ 度量值）。启动链上每一级固件在把控制权交给下一级之前，先把下一级代码的哈希 extend 进 PCR，于是整条启动路径被记录为一组不可伪造的寄存器值。密钥封存（sealing）利用这一点：全盘加密的卷密钥可以按“仅当 PCR 匹配预期值才释放”的条件封存，启动链任何一环被替换，密钥就取不出来。软件入口是 \code{/dev/tpm0} 与操作系统的安全栈（BitLocker、LUKS+TPM、安全启动）。

\begin{longtable}{L{0.15\linewidth}L{0.24\linewidth}L{0.24\linewidth}L{0.30\linewidth}}
\toprule
器件 & 挂接与供电 & 核心职能 & 软件入口 \\
\midrule
RTC & 32.768 kHz 晶振 + 纽扣电池 & 墙钟保持、闹钟/唤醒 & 内核读初值，hwclock，NTP 校准 \\
GPIO & PCH 寄存器块，常电/主电兼有 & 单比特信号进出 & 驱动直接读写，libgpiod \\
SMBus/I2C & PCH 主控，100 kHz & SPD、传感器 & BIOS 读 SPD；i2c 子系统，lm-sensors \\
Super I/O & LPC，主电域 & 串并口、风扇、监控 & hwmon 驱动 \\
EC & LPC/eSPI，常电域 & 电池、风扇、键盘、电源键 & ACPI 方法 + SCI 中断 \\
TPM & LPC/SPI/I2C 或 fTPM & 密钥存储、启动度量 & \code{/dev/tpm0}，全盘加密封存 \\
\bottomrule
\end{longtable}

至此，本节的对象可以收成一句话：整机不是一颗 CPU 加一堆外设，而是一套\hwkey{分层的控制结构}——最快的路径收进 CPU die，慢的长尾留给 PCH 与一众小控制器，层与层之间全部用地址空间中的事务说话。设备枚举决定了每块硬件在地址空间里的位置，中断路由决定了事件如何找到正确的核，IOMMU 决定了设备的 DMA 能碰到哪些内存，定时器决定了操作系统何时醒来；这四件事加上小控制器们的值守，共同构成了 CPU 控制整台机器的控制平面。
% ch06b part3: 第四节 主板：供电、时钟与固件

\section{四、主板：供电、时钟与固件}

前面三节讨论了 CPU 如何通过控制器驾驭内存与外设：命令如何调度、链路如何训练、数据如何在总线上可靠移动。这些芯片不会悬浮在空中工作——它们焊在同一张印制电路板上，由这张板供给能量、供给时钟、供给彼此相连的走线。本节把视角从芯片移到板级：主板这张多层 PCB 究竟提供什么，CPU 核心那约 1V、上百安培的供电由谁稳出来，全机的时钟从哪一颗晶体出发，以及按下电源键之后、操作系统出现之前那段由固件主导的时间。

\topic{4.1\quad 主板是什么：一张板承载三种分配网络}

\term{主板}{motherboard}是一块多层印制电路板（PCB），上面焊接 CPU 插槽、内存插槽、芯片组与各种连接器。从功能上看，主板是三种分配网络的载体。第一种是\hwkey{电源分配网络}（Power Delivery Network，PDN）：24-pin 与 8-pin 电源接口把电源供应器（PSU）的 12V、5V、3.3V 引入板内，VRM 模组（电压调节模组，4.2 节详解）把 12V 降到 CPU 核心需要的约 1V，整面的电源平面与地平面再把电流低阻地送到每个芯片。第二种是\hwkey{时钟分配网络}：一颗 25MHz 晶振驱动时钟发生器芯片，倍频、缓冲、扇出，把参考时钟送到 CPU、内存与每条总线。第三种是\hwkey{信号分配网络}：成千上万的走线按阻抗受控的方式连接各芯片，高速信号贴着完整的参考平面走，长度按组匹配。

三种网络各有自己的约束语言。电源网络的约束是阻抗：负载电流怎么跳变，电压都不能超出允许窗口，这个问题在电源时钟章第一节已用目标阻抗的形式化方法讲过。时钟网络的约束是对齐与纯净：各接收端的沿要齐、沿的位置不能抖动，见该章第二节的时钟树与抖动分析。信号网络的约束是完整：反射、串扰、回流都要受控，见该章第三、四节。主板设计在很大程度上就是这三套约束在同一块层叠结构上的联合求解。

\begin{keyidea}
主板不是被动的连线板。它是三个物理分配网络的交叉点：电流在其中流动，时钟沿在其中传播，信号波在其中往返。CPU 的性能与稳定性，最终受限于这张板把能量、时间和信息送得有多准。
\end{keyidea}

\topic{层叠结构：每一层都有分工}

多层板把铜箔与绝缘介质（FR-4 玻纤环氧，介电常数约 4.2）交替叠合、层压为一整块。层数不是越多越好，而是由信号密度与参考平面需求决定。典型叠层安排如下：

\begin{longtable}{L{0.10\linewidth}L{0.44\linewidth}L{0.20\linewidth}L{0.17\linewidth}}
\toprule
方案 & 各层分工（顶到底） & 适用 & 代价 \\
\midrule
4 层 & 信号/元件、地平面、电源平面、信号 & 低速板、微控制器板 & 只有两层可走线，高速信号密度受限 \\
6 层 & 信号、地、信号、信号、电源/地、信号 & 主流台式机主板 & 中间两层信号共享参考平面，需避免长距离平行走线 \\
8 层 & 信号、地、信号、地、电源、信号、地、信号 & 多通道 DDR、多插槽平台 & 每个高速信号层都有紧邻的完整参考平面，成本显著上升 \\
\bottomrule
\end{longtable}

叠层的核心原则只有一条：高速信号层必须紧邻完整的参考平面（通常是地）。原因在电源时钟章第四节讲过——信号电流沿走线前进，回流沿正下方的平面返回，平面完整则环路面积最小、阻抗连续；平面被割裂，回流被迫绕路，串扰与辐射随之恶化。内存条插槽附近的区域走线最密，DDR 主板因此普遍采用 6 层或 8 层。

\topic{长度匹配：蛇形线是买时间的走线}

同一个数据组内的信号必须同时到达。DDR 的一个字节组里，8 根 DQ 与 1 对 DQS 的走线长度差要控制在零点几毫米以内——FR-4 中信号传播速度约 15cm/ns，即 1ps 对应约 0.15mm，0.3mm 的长度差约折合 2ps 的到达时间差。存储器件结构章第三节讲过，训练可以把每根 DQ 的采样延迟独立调节到数据眼中央，但训练的调节范围有限，板级等长先把粗偏差消掉，训练再做精调，两级分工。走线长度不够时，布线工具把它绕成波浪形的\hwkey{蛇形线}（serpentine）补足长度；蛇形的弯折会局部破坏与邻线的间距，弯折幅度和间距同样有规则约束。差分对（PCIe、USB、SATA）则要求两根线自始至终等长并行，对内失配直接损耗共模抑制能力。

\begin{pdfdiagram}[x=1cm,y=1cm]
  % ---- 电源入口与 VRM（左上） ----
  \node[hw block, minimum width=1.3cm, minimum height=0.8cm, align=center] (p8) at (0.7,8.2) {8-pin\\CPU 供电};
  \node[hw compute, minimum width=2.2cm, minimum height=1.0cm, align=center] (vrm) at (3.1,8.2) {VRM 模组\\多相 buck};
  \node[hw state, minimum width=3.0cm, minimum height=1.7cm, align=center] (cpu) at (6.3,7.7) {CPU 插槽};
  \draw[hw danger] (p8.east) -- (vrm.west);
  \node[hw label] at (1.67,8.48) {12 V};
  \draw[hw danger] (vrm.east) |- (cpu.west);
  \node[hw label] at (4.5,7.88) {VCC};

  % ---- DIMM 插槽（右上） ----
  \node[hw memory, minimum width=0.7cm, minimum height=2.4cm] (dimm1) at (9.3,8.0) {\rotatebox{90}{DIMM}};
  \node[hw memory, minimum width=0.7cm, minimum height=2.4cm] (dimm2) at (10.4,8.0) {\rotatebox{90}{DIMM}};
  \node[hw label] at (9.85,9.5) {DIMM 插槽 $\times$2};
  \draw[hw bus] (cpu.east |- 7.8,7.3) -- (dimm1.west |- 7.8,7.3);
  \draw[hw bus] (dimm1.east |- 7.8,7.3) -- (dimm2.west |- 7.8,7.3);
  \node[hw label, align=center] at (8.37,7.72) {内存\\通道};

  % ---- 24-pin 主供电（右缘） ----
  \node[hw block, minimum width=1.2cm, minimum height=1.2cm, align=center] (p24) at (11.2,6.2) {24-pin\\主供电};
  \draw[BookRed, line width=0.62pt] (p24.west) -- (9.7,6.2);
  \node[hw label, font=\tiny] at (10.1,5.82) {12 V/5 V/3.3 V 入板};

  % ---- PCIe 插槽（左中、左下） ----
  \node[hw block, minimum width=4.8cm, minimum height=0.7cm] (pcie16) at (3.0,5.4) {PCIe $\times$16 插槽};
  \node[hw block, minimum width=3.6cm, minimum height=0.6cm] (pcie1) at (2.4,3.7) {PCIe $\times$1 插槽};
  \draw[hw bus] (5.0,6.85) -- (5.0,5.75);

  % ---- PCH 与低速芯片 ----
  \node[hw control, minimum width=2.2cm, minimum height=1.1cm] (pch) at (6.6,4.6) {PCH};
  \draw[hw bus] (6.9,6.85) -- (6.9,5.15);
  \node[hw label, anchor=east, fill=white, inner sep=1pt] at (6.75,5.95) {DMI};
  \draw[hw bus] (pch.west |- 5.5,4.4) -- (4.8,4.4) |- (4.8,3.7) -- (4.2,3.7);
  \node[hw label, anchor=east] at (4.7,4.2) {PCIe $\times$1};
  \node[hw memory, minimum width=1.7cm, minimum height=0.8cm, align=center] (spi) at (7.4,3.0) {SPI flash\\固件};
  \node[hw compute, minimum width=2.0cm, minimum height=0.8cm] (sio) at (9.9,3.0) {Super I/O};
  \draw[hw bus] (6.2,4.05) |- (6.2,3.0) -- (spi.west);
  \draw[hw bus] (spi.east) -- (sio.west);
  \node[hw label, anchor=east, fill=white, inner sep=1pt] at (6.05,3.55) {SPI / eSPI};

  % ---- 时钟发生器与扇出 ----
  \node[hw control, minimum width=2.0cm, minimum height=0.9cm, align=center] (clk) at (9.4,4.8) {时钟\\发生器};
  \draw[hw return] (clk.north) -- (9.4,6.8);
  \draw[hw return] (9.4,6.5) -- (7.4,6.5) -- (7.4,6.85);
  \draw[hw return] (9.4,6.05) -- (2.4,6.05) -- (2.4,5.75);
  \draw[hw return] (9.4,5.6) -- (8.2,5.6) -- (8.2,4.6) -- (7.7,4.6);
  \node[hw label, anchor=south] at (4.0,6.14) {100 MHz};

  % ---- 图例 ----
  \draw[hw bus] (0.6,0.75) -- (1.7,0.75);
  \node[hw label, anchor=west] at (1.9,0.75) {信号走线（数据/地址/控制）};
  \draw[hw danger] (0.6,0.3) -- (1.7,0.3);
  \node[hw label, anchor=west] at (1.9,0.3) {电源分配（12 V 输入、VRM 输出）};
  \draw[hw return] (6.2,0.75) -- (7.3,0.75);
  \node[hw label, anchor=west] at (7.5,0.75) {时钟分配（100 MHz 参考）};
\end{pdfdiagram}
\diagramnote{台式机主板布局与三大分配网络。实线为信号路径：CPU 直连内存与 PCIe $\times$16，其余设备经 DMI 汇聚到 PCH。红色线为电源路径：8-pin 专供 CPU，经 VRM 降为约 1V；24-pin 把 12V/5V/3.3V 引入全板。虚线为时钟路径：时钟发生器向 CPU、PCH、PCIe 扇出 100MHz 参考。SPI flash 保存固件，Super I/O 承接键盘、串口、风扇监控等低速功能。}

读这张图时值得注意谁在 CPU 上、谁在 PCH 上：时延最敏感的内存与显卡直连 CPU，其余一切低速路径汇聚到 PCH，再经一条 DMI 链路上行——这正是经典芯片与平台章讲过的南北桥职责划分的现代形态。主板上还能看到一颗 Super I/O 芯片（职能见 3.6 节），它经 LPC 或 eSPI（Enhanced Serial Peripheral Interface，增强型串行外设接口）挂在 PCH 上，承接比 PCH 更慢、更杂的功能。

\topic{4.2\quad VRM：一个模拟反馈环给 CPU 供电}

CPU 核心是主板上最苛刻的负载：电压只有约 1V，电流却可达上百安培，而且电流随负载在纳秒到微秒尺度上剧烈变化。给它供电的电路称为\term{电压调节模组}{Voltage Regulator Module，VRM}，其主体是多相并联的 buck 降压变换器。电源时钟章第一节比较过线性稳压与开关稳压的取舍：压差大、电流大时必须用开关方案。12V 降到 1V、输出上百安培，正是开关方案的主场；本节把这个方案拆开，看它的功率级、反馈环与瞬态行为。

\topic{buck 单相：占空比决定电压}

\term{buck 变换器}{buck converter}的功率级只有四个角色：高边 MOSFET、低边 MOSFET、电感、输出电容。两只 MOSFET 组成半桥，中点称为开关节点 SW。高边导通时 SW 接到 12V，电感电流线性上升，电感储能；高边关断、低边导通时 SW 接地，电感电流经低边续流，线性下降，把储存的能量释放给输出。控制电路以固定频率（每相典型 300kHz--1MHz）交替驱动两只管子，高边导通时间占周期的比例称为占空比 $D$。稳态下电感电压平均值为零，输出电压就是 SW 波形的平均值：

$$V_{out} = D \cdot V_{in}$$

12V 转 1V 意味着 $D \approx 0.083$：每个周期里高边只导通约 8\% 的时间。两只管子绝不允许同时导通——那会把 12V 直接短路到地——驱动器在两次导通之间插入几十纳秒的死区时间。电感电流的三角起伏称为纹波电流，由输出电容滤平，输出电压上只残留很小的锯齿。

\begin{pdfdiagram}[x=1cm,y=1cm]
  % ---- 功率级（上方，实线） ----
  \node[hw block, minimum width=1.3cm, minimum height=0.9cm, align=center] (vin) at (0.75,4.8) {12 V\\输入};
  \node[hw compute, minimum width=2.0cm, minimum height=1.8cm, align=center] (bridge) at (3.0,4.8) {半桥\\高边+低边\\MOSFET};
  \node[hw compute, minimum width=1.4cm, minimum height=0.9cm] (ind) at (5.6,4.8) {电感 L};
  \node[hw state, minimum width=1.8cm, minimum height=1.0cm, align=center] (load) at (8.6,4.8) {负载\\CPU 核心};
  \draw[hw bus] (vin.east) -- (bridge.west);
  \draw[hw bus] (bridge.east) -- (ind.west);
  \node[hw label] at (4.45,5.02) {SW};
  \draw[hw bus] (ind.east) -- (load.west);
  \node[hw label] at (7.0,5.1) {$V_{out}\approx 1$ V};
  \node[hw compute, minimum width=1.5cm, minimum height=0.8cm] (cap) at (6.6,3.6) {输出电容 C};
  \draw[line width=0.62pt, draw=BookMuted] (6.6,4.8) -- (cap.north);

  % ---- 反馈环路（下方，虚线） ----
  \node[hw compute, minimum width=2.2cm, minimum height=0.9cm, align=center] (div) at (5.6,1.5) {分压采样 R1/R2};
  \node[hw compute, minimum width=1.8cm, minimum height=1.0cm, align=center] (ea) at (2.9,1.5) {误差\\放大器};
  \node[hw memory, minimum width=1.3cm, minimum height=0.6cm] (vref) at (2.9,0.3) {$V_{ref}$};
  \node[hw control, minimum width=1.8cm, minimum height=1.0cm, align=center] (pwm) at (2.9,2.9) {PWM\\比较器};
  \node[hw block, minimum width=1.3cm, minimum height=0.8cm] (saw) at (0.85,2.9) {锯齿波};
  \draw[hw return] (load.south) -- (8.6,1.5) -- (div.east);
  \draw[hw return] (div.west) -- (ea.east);
  \draw[hw bus] (vref.north) -- (2.9,1.0);
  \draw[hw return] (ea.north) -- (2.9,2.4);
  \draw[hw bus] (saw.east) -- (pwm.west);
  \draw[hw bus] (pwm.north) -- (2.9,3.9);
  \node[hw label, anchor=west] at (3.05,3.65) {栅极驱动};
\end{pdfdiagram}
\diagramnote{buck 单相功率级与反馈环。实线为功率路径：半桥把 12V 斩成方波，电感与输出电容滤成直流。虚线为反馈路径：输出电压经分压采样送入误差放大器，与参考电压比较，误差信号再与锯齿波比较产生 PWM，决定下一个周期高边的导通时间。输出偏低则占空比加大，输出偏高则占空比减小——一个自动回到设定值的负反馈环。}

\topic{反馈环是一个模拟 PID}

图的下半部分是这个变换器的控制核心，也是本章主题在电源领域的化身。输出电压经电阻分压采样后送入误差放大器的反相端，同相端接精确的参考电压（带隙基准）；误差放大器输出两者之差经补偿网络放大后的信号。这个误差信号送到 PWM 比较器，与固定频率的锯齿波比较：误差大，交点右移，占空比加大；误差小，占空比减小。输出电压任何偏离设定值的趋势都会被环路反向纠正——结构与负反馈放大器完全同构——本章第五节将给出它的统一词汇，只是这里的反馈量不是数字状态而是模拟电压。

误差放大器周围的补偿网络决定环路的动态性格：纯比例增益响应快但留静差，积分环节消除静差，微分环节抑制振荡。补偿元件（几只电阻电容）取值合起来就是一个模拟 PI/PID 控制器，其设计目标是把环路带宽推到开关频率的约十分之一（典型几十到一百多 kHz），同时保留足够的相位裕度。带宽就是这个电源“反应速度”的定量刻画，它直接决定下一段要讲的负载瞬态表现。

现代 VRM 的参考电压不是固定的。CPU 通过一条三线串行总线（Intel 平台称 SVID，即 Serial VID；AMD 平台对应 SVI2/SVI3）向 VRM 控制器发送 \hwkey{VID}（Voltage Identification，电压识别码），以 5--10mV 的步进指定目标电压：轻载时请求降压省电，重载冲刺前请求升压。操作系统层面的频率调节（P-state 切换），对应到硬件上就是一次 VID 改写加一次倍频改写——软件看见的“变频”，有一半是这个模拟环路在换设定值。

\topic{多相并联：错相 360/N 度}

单相 buck 的电流能力受 MOSFET 与电感限制，每相持续输出典型 20--40A。CPU 核心需要上百安培时，唯一的出路是多相并联：$N$ 个相同的半桥加电感共用同一组输出电容，由同一个控制器驱动。台式机主板常见 4--16 相，相数印在主板规格表上，是其供电能力的直接指标。

多相的价值不止于电流相加。控制器让各相的 PWM 依次错开 $360/N$ 度：6 相错 60 度，任一时刻只有一相的高边刚导通。各相电感电流的三角纹波相位错开，在输出电容处汇合时峰谷互抵，合成纹波的频率变为单相的 $N$ 倍、幅度显著减小——同样的滤波电容得到更平的输出。输入侧同样受益：12V 输入电流被摊成 $N$ 份错相抽取，输入电容与 PSU 看到的脉冲电流随之减小。瞬态时多相可以同时转向，等效电感为单相的 $L/N$，电流爬升速度也是单相的 $N$ 倍。

\topic{负载瞬态：下陷与过冲}

buck 环路再快也有带宽上限，而 CPU 的电流阶跃几乎没有上升时间：一批核心从空闲同时进入睿频，电流边沿的 di/dt 可达几十 A/ns，一次完整的负载阶跃（如 10A 跳到 100A）在几十纳秒到微秒量级内完成。电感电流不能突变，环路察觉到电压下跌、调大占空比也需要若干微秒。这段响应延迟里的电流缺口只能由输出电容垫付，电压因此下跌，称为\hwkey{下陷}（droop/undershoot）；反过来，满载突然撤去时电感里已建立的电流无处释放，灌进输出电容造成\hwkey{过冲}（overshoot）。

下陷幅度可以估算。缺口电流由电容提供，$\Delta V \approx \Delta I \cdot \Delta t / C$，再叠上电容等效串联电阻造成的 $\Delta I \cdot \mathrm{ESR}$ 一项。代入典型量级：$\Delta I = 90$A，环路响应延迟 $\Delta t \approx 2\,\mu$s，输出电容总量 $1.5$mF，得到 $\Delta V \approx 90 \times 2\,\mu\mathrm{s} / 1.5\,\mathrm{mF} = 120$mV——对 1V 供电已是 12\% 的扰动，而 CPU 要求的供电窗口通常只有正负几个百分点。这就是主板上 CPU 插槽周围密密麻麻排布电容的原因：聚合物电容管低频段的大电荷量，成片的 MLCC 陶瓷电容管高频段的快速响应，分工方式与电源时钟章第一节的去耦分层完全相同。

\begin{longtable}{L{0.30\linewidth}L{0.24\linewidth}L{0.38\linewidth}}
\toprule
量 & 典型量级 & 由什么决定 \\
\midrule
负载电流阶跃 $\Delta I$ & 数十到上百 A & 核心数、睿频策略、负载类型 \\
电流变化率 di/dt & 边沿可达几十 A/ns & CPU 内部时钟门控的速度，主板无法控制 \\
环路响应延迟 & 数 $\mu$s & 环路带宽（开关频率的约 1/10）与补偿设计 \\
下陷/过冲幅度 & 数十到一百多 mV & $\Delta I$、响应延迟、输出电容总量与 ESR \\
供电允许窗口 & 约 $\pm$3\%--5\% & CPU 数据手册的 VRM 规范 \\
\bottomrule
\end{longtable}

\topic{负载线：故意不完美的稳压}

直觉上稳压器的目标是输出电压恒定，VRM 却故意做不到：它让输出电压随负载电流线性下降，满载时比空载低几十到一百多 mV。这条人为的伏安斜线称为\term{负载线}{load line}，其斜率（等效电阻）典型为 1m$\Omega$ 量级，即每 100A 负载压降约 100mV。

故意降压的理由是双向余量。没有负载线时，空载 1.00V、加载下陷 120mV，电压触及 0.88V，越出下限；有了负载线，满载稳态值预先设在 0.90V 附近，同样的下陷从更低处开始、谷底不变，而撤载过冲从 0.90V 起跳，顶点远低于从 1.00V 起跳的情形——同一份电压窗口，瞬态扰动在上下两个方向都获得了空间。附带的收益是省电：满载时电压低 10\%，同样的电流下功耗也低约 10\%（理想化估算——实际负载电流会随电压略变）。负载线的斜率由 VRM 控制器按 CPU 厂商规范实现，超频工具里的 LLC（Load-Line Calibration）档位调的就是这条线的斜率——把线拉平，满载电压更好看，代价是撤载过冲顶到上限。

\begin{warningbox}
“供电相数越多电压越稳”只说对了一半。相数决定电流能力与纹波抵消，瞬态表现却由环路带宽、输出电容与负载线共同决定；两相供电配优秀的补偿与充足的电容，瞬态响应可以好过堆相数但补偿平庸的设计。看一份 VRM 方案，相数、每相电流等级、控制器带宽、输出电容总量要放在一起读。
\end{warningbox}

\topic{4.3\quad 时钟分配：一颗晶体驱动全机}

主板上的时钟起点是一颗 25MHz 石英晶体（或晶体振荡器），它本身只产生一个频率。晶体旁边是\term{时钟发生器}{clock generator}芯片：内部 PLL 把 25MHz 倍频到 100MHz，再经一组扇出缓冲器复制成多路，分别送往 CPU（作为基础时钟 BCLK）、PCH、PCIe 插槽与 SATA 等接口。PLL 的倍频与滤波原理、晶振起振需要毫秒级时间的事实，电源时钟章第二节都已讲过，此处不再重复；本节关心的是分配的结构。

各域拿到的是同一份 100MHz 参考，各自的内部 PLL 再做二次倍频：CPU 核心把 BCLK 乘上倍频系数（如 $\times 45$ 得 4.5GHz），内存控制器把参考倍频到内存实际时钟（DDR4-3200 对应 1600MHz），PCIe 与 SATA 则直接用 100MHz 参考恢复出各自的链路速率。同源设计的好处是各域时钟的漂移方向一致，跨域接口看到的频率关系是稳定的。

\begin{pdfdiagram}[x=1cm,y=1cm]
  \node[hw block, minimum width=1.4cm, minimum height=0.8cm, align=center] (xtal) at (0.9,3.0) {25 MHz\\晶振};
  \node[hw control, minimum width=2.6cm, minimum height=1.2cm, align=center] (gen) at (4.0,3.0) {时钟发生器\\PLL + 扇出缓冲};
  \node[hw state, minimum width=2.6cm, minimum height=0.9cm, align=center] (cpu) at (9.3,5.0) {CPU 核心\\$\times$倍频系数};
  \node[hw compute, minimum width=2.6cm, minimum height=0.9cm, align=center] (mem) at (9.3,3.6) {内存控制器\\$\times$16 等};
  \node[hw compute, minimum width=2.6cm, minimum height=0.9cm, align=center] (io) at (9.3,2.2) {PCIe / SATA\\直接用参考};
  \draw[hw bus] (xtal.east) -- (gen.west);
  \draw[hw bus] (gen.east) -- (6.3,3.0);
  \draw[hw bus] (6.3,3.0) -- (6.3,5.0) -- (cpu.west);
  \draw[hw bus] (6.3,3.0) -- (6.3,3.6) -- (mem.west);
  \draw[hw bus] (6.3,3.0) -- (6.3,2.2) -- (io.west);
  \fill[BookMuted] (6.3,3.6) circle (0.045);
  \fill[BookMuted] (6.3,2.2) circle (0.045);
  \node[hw label] at (7.15,5.26) {100 MHz 参考};
  \node[hw label, align=center] at (4.0,1.7) {扩频（SSC）：约 33 kHz 三角波，频率在\\标称值与其 $-0.5\%$ 之间扫动（下扩频）};
\end{pdfdiagram}
\diagramnote{主板时钟分配。25MHz 晶振驱动时钟发生器，PLL 倍频出 100MHz 参考并扇出到各时钟域；各域内部 PLL 二次倍频。SSC 调制在时钟发生器内完成，全机时钟同步漂移。}

\topic{扩频时钟：把能量摊开来过认证}

恒频时钟的频谱是一根尖锐的谱线：100MHz 及其各次谐波上的能量高度集中，主板上长走线充当天线时，这些频点的辐射容易超出电磁兼容限值。\term{扩频时钟}{Spread Spectrum Clocking，SSC}的做法是让时钟频率不再恒定，而是以约 33 kHz 的三角波缓慢调制，使频率在标称值与标称值的 $-0.5\%$ 之间来回扫动（下扩频）。总能量不变，但原本集中在一个频点的能量被摊进约 0.5\% 带宽内，峰值功率谱密度随之下降数 dB 到十 dB 量级，恰好够把认证测试中的辐射峰值降到限值以下。代价是时钟沿位置引入了有界的低频摆动：平均频率比标称值低 0.25\%，瞬时频率在标称值与其 $-0.5\%$ 之间变化。同步数字逻辑对此不敏感；高速串行链路则要求收发双方共享同一份带 SSC 的参考——PCIe 允许 SSC 的前提正是主板把同一路 100MHz 参考同时发给 CPU 根复合体与插槽，两者同步漂移，相对关系不变。电磁兼容的完整讨论见电源时钟章第四节。

\topic{4.4\quad 固件与开机：从 PWR\_OK 到 bootloader}

按下电源键到操作系统出现，中间有一段完全由固件主导的时间。经典芯片与平台章已讲过一个版本的故事：复位向量、BIOS 的 POST、设备初始化、选设备与移交四阶段。本节讲现代 UEFI 版本，重点放在这条链路与主板硬件的对应关系——固件放在哪颗芯片里，每一步初始化的是本节与前几节讲过的哪部分硬件。

\topic{上电时序：先电力，再时钟，最后逻辑}

电源键并不直接接通主电源，它只是给 PCH（或嵌入式控制器 EC）一个请求信号。PCH 确认后把电源控制线 PS\_ON\# 拉低（低有效），PSU 才开始输出各电压轨；板上各 VRM 随之启动，把 12V 分配到 CPU、内存与各插槽。电压爬升与稳定需要时间，PSU 在所有输出进入允许范围后才拉高 PWR\_OK 信号，\hwevidence{ATX 规范把这个延迟规定在 100--500ms}。与此同时晶振起振（毫秒级）、时钟发生器内 PLL 锁定（数十到数百 $\mu$s，见电源时钟章第二节）。PCH 等到 PWR\_OK 有效且时钟稳定，才释放全机复位。CPU 退出复位后的第一个动作是从复位向量取指——x86 的初始取指地址在 4GB 空间顶端附近（\code{0xFFFFFFF0}），这个地址由 PCH 译码，映射到主板上的 SPI flash。复位向量在不同架构上的对照，经典芯片与平台章已有专表，此处只需记住结论：入口地址是硬件约定，响应该地址的必须是可靠的非易失存储。

\topic{UEFI 的阶段：SEC、PEI、DXE}

固件保存在主板上那颗 SPI NOR flash 里，典型容量 16--32MB，经 PCH 的 SPI 控制器以几十 MHz 时钟读取，并映射到地址空间顶端供 CPU 直接取指。UEFI 固件把启动过程组织成几个阶段，每个阶段的可用资源不同，能做的事也不同：

\begin{longtable}{L{0.09\linewidth}L{0.24\linewidth}L{0.33\linewidth}L{0.25\linewidth}}
\toprule
阶段 & 运行环境 & 主要工作 & 结束时的世界 \\
\midrule
SEC & 无内存可用，单核，实模式入口，随即切保护模式 & 把部分 cache 配置为临时 RAM（cache-as-RAM），建立栈与初始数据区，验证并装载 PEI 核心 & CPU 有了可用的临时内存 \\
PEI & 临时 RAM，SPI flash 可读 & 早期硬件初始化；执行内存参考代码：读 SPD、配置内存控制器、执行训练；内存就绪后装载 DXE 核心 & DRAM 可用 \\
DXE & 全部内存可用 & 按依赖关系调度驱动：枚举 PCIe 设备、初始化存储与显示控制器、建立 ACPI/SMBIOS 表 & 设备可枚举、表已备好 \\
BDS & 设备驱动就绪 & 初始化控制台，按启动项顺序选择启动设备，装载 bootloader & 控制权移交 \\
\bottomrule
\end{longtable}

SEC 阶段的 cache-as-RAM 值得停留片刻。复位后 DRAM 尚未初始化——它需要先读 SPD、再做训练（本章 2.3 节），而做这些事的代码本身需要内存来运行。解决方法是把 CPU 的部分 cache 临时配置成填充后不再驱逐、也不写回（no-evict）的模式，让它充当几十 KB 到几 MB 量级的 SRAM：C 代码的栈和全局变量就建在这片临时内存上。等到 PEI 阶段末尾 DRAM 训练完成，固件把执行环境整体迁入真正的内存，cache 才恢复本来的职责。PEI 阶段的内存初始化是整机启动中最容易失败的一环：训练失败时，主板诊断灯停在 DRAM 档位，正是这个阶段的输出。

\begin{pdfdiagram}[x=1cm,y=1cm]
  \node[hw block, minimum width=6.0cm, minimum height=0.85cm] (b1) at (5.5,9.6) {按下电源键：PSU 各电压轨升起};
  \node[hw control, minimum width=6.0cm, minimum height=0.85cm] (b2) at (5.5,8.4) {PWR\_OK 拉高，复位释放};
  \node[hw state, minimum width=6.0cm, minimum height=0.85cm] (b3) at (5.5,7.2) {CPU 从复位向量取指（命中 SPI flash）};
  \node[hw compute, minimum width=6.0cm, minimum height=0.85cm] (b4) at (5.5,6.0) {SEC：cache-as-RAM，建立临时执行环境};
  \node[hw compute, minimum width=6.0cm, minimum height=0.85cm] (b5) at (5.5,4.8) {PEI：读 SPD，内存初始化与训练};
  \node[hw compute, minimum width=6.0cm, minimum height=0.85cm] (b6) at (5.5,3.6) {DXE：枚举设备，装载驱动，建立 ACPI 表};
  \node[hw state, minimum width=6.0cm, minimum height=0.85cm] (b7) at (5.5,2.4) {BDS：选择启动设备，装载 bootloader};
  \node[hw block, minimum width=6.0cm, minimum height=0.85cm] (b8) at (5.5,1.2) {交棒：bootloader $\to$ OS 内核};
  \draw[hw bus] (b1.south) -- (b2.north);
  \draw[hw bus] (b2.south) -- (b3.north);
  \draw[hw bus] (b3.south) -- (b4.north);
  \draw[hw bus] (b4.south) -- (b5.north);
  \draw[hw bus] (b5.south) -- (b6.north);
  \draw[hw bus] (b6.south) -- (b7.north);
  \draw[hw bus] (b7.south) -- (b8.north);
  \node[hw label, anchor=east] at (2.3,8.4) {等电压与时钟稳定};
  \node[hw label, anchor=east] at (2.3,6.0) {DRAM 尚不可用};
  \node[hw label, anchor=east] at (2.3,4.8) {SPD 见 4.5 节};
  \node[hw label, anchor=west, align=left] at (8.7,7.2) {复位向量对照见\\经典芯片与平台演进章};
  \node[hw label, anchor=west] at (8.7,3.6) {表留给 OS 解析};
\end{pdfdiagram}
\diagramnote{UEFI 开机流程。前半段被硬件约束主导：电压与时钟稳定之前不能释放复位，内存训练完成之前只能用 cache 充当内存；后半段被软件结构主导：驱动按依赖调度，表建好后才交棒。任何一步失败，整机表现为“点不亮”，而根因可能远在 CPU 之外。}

\topic{ACPI 表：固件写给操作系统的自我介绍}

bootloader 与操作系统接手后，面对的第一个问题是：这台机器有几个核、中断控制器在哪、有哪些设备？现代平台不允许 OS 靠试探端口来回答——太危险也太慢。答案是 \term{ACPI}{Advanced Configuration and Power Interface，高级配置与电源接口}：固件在 DXE 阶段把一组表写进内存，并在约定位置留下签名指针 \code{RSDP}；OS 由指针找到根表 \code{XSDT}，再按其中的条目索引到各张子表。硬件清单以数据而非探测的方式交付。

\begin{longtable}{L{0.14\linewidth}L{0.40\linewidth}L{0.38\linewidth}}
\toprule
表 & 内容 & OS 拿它做什么 \\
\midrule
MADT & 每个处理器核的条目、本地 APIC 地址（典型 \code{0xFEE00000}）、I/O APIC 与中断源覆盖 & 决定唤醒几个核、中断控制器在哪、中断号如何路由 \\
MCFG & PCIe 增强配置空间（ECAM）的基址 & 用 MMIO 读写所有 PCIe 设备的配置空间（见主存与总线章的 MMIO） \\
DSDT & Differentiated System Description Table（系统描述表）：AML 字节码描述的设备树；设备对象、\code{_HID} 即插即用 ID、\code{_CRS} 资源声明 & OS 的 ACPI 解释器执行 AML，枚举无标准总线的设备（电源键、温度传感器等） \\
FADT & 电源管理相关的固定寄存器块地址与系统属性 & 睡眠/唤醒、电源按钮语义、复位寄存器位置 \\
\bottomrule
\end{longtable}

这套机制把“主板厂商才知道的事”——哪个引脚接了哪个设备、电源键走哪条路径——封装成标准数据结构。同一份 OS 镜像因此能在不同主板上启动：差异被压缩在表的内容里，而不是代码里。BDS 阶段的交棒本身很简短：固件从 EFI 系统分区（ESP）读取启动项指定的 \code{.efi} 文件（bootloader），把它装入内存并跳转；此后固件退居幕后，只保留 ACPI 表与少量运行时服务。从 BIOS 四阶段到 UEFI 的 SEC/PEI/DXE/BDS，载体换了，骨架未变——经典芯片与平台章末尾的四阶段对照依然成立。

\topic{4.5\quad SPD：内存条身上的自我介绍}

PEI 阶段配置内存控制器时，时序参数从哪来？内存条自己带着答案。每根 DIMM 上都有一颗小 EEPROM，即 3.6 节介绍过的 SPD（Serial Presence Detect，串行存在检测），\hwevidence{DDR4 的 SPD 容量为 512 字节（JEDEC EE1004 规范）}；DDR5 把 SPD 换成 1024 字节的 SPD5118 集线器，接口也换成 I3C，角色不变。整颗芯片里固化着这根条的完整身份与能力档案。主板通过 I2C/SMBus 总线（7 位地址 \code{0x50}--\code{0x57}，标准模式 100kHz）在上电早期把它读出来——读满 512 字节只需几十毫秒。

\begin{longtable}{L{0.22\linewidth}L{0.42\linewidth}L{0.28\linewidth}}
\toprule
字段 & 内容 & 用途 \\
\midrule
制造信息 & 厂商、颗粒型号、序列号、生产日期 & 固件与诊断工具识别条子身份 \\
组织结构 & 容量、rank 数、位宽、bank 数 & 内存控制器配置地址映射 \\
标准时序表 & 各档频率对应的最小 \hwtiming{tCK、CL、tRCD、tRP、tRAS}（如 DDR4-3200：\hwtiming{tCK 0.625ns、CL 22}） & 决定目标频率下的命令间隔 \\
标称电压 & 标准档 1.2V，性能档所需电压 & VRM 设定内存供电 \\
扩展 profile & Intel XMP / AMD EXPO：厂商验证过的高频时序组合 & 用户在固件设置里一键启用超频档 \\
\bottomrule
\end{longtable}

SPD 把本章的两条线索接在了一起。4.4 节讲过，PEI 阶段的内存参考代码负责点亮 DRAM；本章 2.3 节讲过，点亮的关键是训练——把每条链路的采样延迟与判决门限调到数据眼中央。SPD 正是这两者的衔接物：固件先按 SPD 中的保守参数（或标称档）让链路以安全时序跑起来，使基本读写可用；训练再以这个工作点为起点做精调，最终把各档位时序写进控制器寄存器。存储器件结构章第三节描述的训练序列因此有一个隐含的前提：训练开始前，控制器已经通过 SPD 知道了这根条子的组织结构和目标时序。没有 SPD，主板无法区分一根 8GB 单 rank 的 3200 条和一根 32GB 双 rank 的 4800 条，也就无从谈起正确的初始化。

SPD 也是“软件怎么碰到硬件”的一个干净样本：它不产生数据流，不搬指令，只是一份躺在 EEPROM 里的参数表，却决定了训练从哪个工作点开始、内存最终跑在什么频率上。用户在固件设置界面里打开 XMP，效果就是把 SPD 扩展区里那组更激进的时序换成训练的目标值——一次配置选择，穿过固件、控制器与 PHY，最终改变的是纳秒级的信号采样位置。

\begin{classicnote}
把本节与前几章连起来读：VRM 的反馈环（4.2）是负反馈放大器结构的功率版本；时钟分配（4.3）是时钟树在板级的延伸；SEC 的 cache-as-RAM（4.4）是最小 CPU 章“取指需要可用的存储”这一约束在真实启动中的解法；SPD 与训练（4.5）则完成了从一颗 EEPROM 到纳秒级采样位置的完整链路。整机的启动不是某个部件的动作，而是这些机制按依赖顺序依次成立的过程。
\end{classicnote}
\section{五、机器里的控制理论：反馈环、稳定性与 PID}

到这里，本章已经出现了好几个控制环：VRM 把输出电压锁在参考值上，PLL 把时钟相位锁在参考沿上，固件把芯片温度维持在温度预算内、把功耗维持在功耗预算内，风扇按温度调节转速，内存训练把采样点定在眼图中心。它们的物理量、执行机构、响应速度完全不同——快的以纳秒计，慢的以秒计——但抽象成框图之后，是同一个数学结构：测量实际值，与目标比较，按差值施加纠正。这一节把这个结构本身讲清楚：先建立反馈控制的通用词汇，再回答一个任何控制环都绕不开的问题——它会不会振荡，最后用四个案例把理论读回机器里的具体电路。掌握这一节后，再遇到数据手册里"环路带宽""补偿网络""锁定时间"这类字样，读到的将不再是孤立的参数，而是同一个结构在不同物理载体上的表现。

\begin{keyidea}
机器里凡是要"把某个物理量维持在目标值"的地方，都有一个反馈环：传感器测量实际值，控制器按误差计算纠正量，执行器把纠正量施加到被控对象上。反馈的代价是稳定性问题——纠正信号在环路里绕行一周带回的滞后，可能让"纠正"在某个频率上变成"助推"。控制设计的全部手艺，就是在响应速度与稳定裕量之间做取舍。
\end{keyidea}

\topic{5.1\quad 反馈控制的基本词汇}

控制理论有一套通用词汇，每个词在机器里都能找到对应的实物。先用机器里的例子把它们逐个定义清楚。

\term{被控对象}{plant}是要被维持状态的物理系统，它的某个物理量是被控量。VRM 电压环里，被控对象是 buck 变换器的 LC 输出级连同挂在上面的负载，被控量是输出电压；PLL 里，被控对象是压控振荡器，被控量是输出时钟的相位；HDD 伺服里，被控对象是磁头臂的机械结构，被控量是磁头相对磁道中心的位置；热管理环里，被控对象是硅片与散热器组成的热系统，被控量是芯片结温。

\term{传感器}{sensor}把被控量转换成可以与目标比较的电信号。VRM 用两只精密电阻组成的分压网络把输出电压按比例缩小；PLL 用鉴相器把两个时钟沿的先后变成脉冲宽度；CPU 用集成在硅片上的热敏二极管或数字温度传感器报告结温；HDD 伺服从磁盘上的伺服扇区解调出位置误差信号（存储器件结构章第六节）。

\term{执行器}{actuator}是控制器能够推动的物理入口。buck 功率级的占空比、VCO 的控制电压、音圈电机的驱动电流、风扇的 PWM 占空比，都是执行器的输入。控制器的全部影响力都通过执行器进入物理世界，执行器的速度上限也因此决定了整个环路的响应上限。

\term{参考值}{reference, setpoint}是目标本身。VRM 的参考多由带隙基准派生——带隙基准是一个对温度和工艺都高度稳定的片上电压源（约 1.2 V），从它派生的参考电压典型在 0.6--1.25 V，经分压比折算后对应目标输出电压；PLL 的参考是外部晶体提供的时钟沿；热环的参考是写入寄存器的温度阈值。

\term{误差}{error}是参考值与测量值之差，$e = r - \hat{y}$。控制器的输入只有这一个数：它不直接知道被控对象发生了什么，只知道"偏了多少、往哪边偏"。理解反馈环的一个好习惯是追问：这个环的误差信号到底是什么物理形式？VRM 里是误差放大器两输入端的电压差，PLL 里是鉴相器输出脉冲的宽度，HDD 伺服里是 PES 数值。

最后一个区分是\term{前馈}{feedforward}与\term{反馈}{feedback}。反馈测量\emph{结果}再纠正：误差出现之后才行动，因此它能抑制模型里没有预料到的扰动，代价是慢——至少要等扰动造成可测量的偏差。前馈测量\emph{扰动}提前补偿：知道负载电流变了多少，就预先多给一份占空比；知道盘片每转的偏心规律，就按查表值预先修正磁头位置（HDD 的 RRO——repeatable runout，重复性跑偏——前馈表，见存储器件结构章第六节）。前馈不等误差出现，因此快，但它只能补偿模型里写明的扰动，模型之外的一概无效。工程上的成熟做法是让两者分工：前馈处理已知的、可预测的扰动，反馈收拾其余一切。

\begin{pdfdiagram}[x=1cm,y=1cm]
  % 主通路：参考 -> 求和点 -> 控制器 -> 执行器 -> 被控对象 -> 输出
  \node[hw label, anchor=east] at (-0.1,0) {参考值 $r$};
  \draw[hw bus] (0.0,0) -- (0.88,0);
  \draw[BookMuted, line width=0.6pt] (1.2,0) circle (0.3);
  \node[hw label, font=\tiny] at (0.84,0.16) {$+$};
  \node[hw label, font=\tiny, anchor=west] at (1.32,-0.38) {$-$};
  \node[hw control, minimum width=1.5cm] (ctrl) at (2.9,0) {控制器};
  \node[hw compute, minimum width=1.5cm] (act) at (5.1,0) {执行器};
  \node[hw compute, minimum width=1.6cm] (plant) at (7.3,0) {被控对象};
  \draw[hw bus] (1.5,0) -- (ctrl.west);
  \node[hw label, font=\tiny] at (1.85,0.22) {误差 $e$};
  \draw[hw bus] (ctrl.east) -- (act.west);
  \draw[hw bus] (act.east) -- (plant.west);
  \draw[hw bus] (plant.east) -- (9.5,0);
  \node[hw label, anchor=west] at (9.6,0) {输出 $y$};
  \fill[BookMuted] (9.0,0) circle (0.05);
  % 扰动从上方进入被控对象
  \draw[hw danger] (7.3,1.0) -- (plant.north);
  \node[hw label, font=\tiny] at (7.3,1.28) {扰动：负载突变、温度、噪声};
  % 反馈通路：输出 -> 下 -> 传感器 -> 左 -> 上进求和点
  \node[hw state, minimum width=1.5cm] (sens) at (5.1,-1.9) {传感器};
  \draw[hw return] (9.0,0) -- (9.0,-1.9) -- (sens.east);
  \draw[hw return] (sens.west) -- (1.2,-1.9) -- (1.2,-0.32);
  \node[hw label, font=\tiny, fill=white, inner sep=1pt] at (3.0,-2.06) {测量值 $\hat{y}$};
  % 机器里的实例标注（各框下方）
  \node[hw label, font=\tiny, text width=2.6cm] at (2.9,-0.85) {误差放大器 / 环路滤波 / 伺服 PID};
  \node[hw label, font=\tiny, text width=2.4cm] at (5.1,-0.85) {buck 功率级 / VCO / 音圈电机};
  \node[hw label, font=\tiny, text width=2.4cm] at (7.3,-0.85) {LC 输出与负载 / 磁头机械结构};
  \node[hw label, font=\tiny, text width=2.4cm] at (5.1,-2.75) {分压网络 / 鉴相器 / 温度传感器};
\end{pdfdiagram}
\diagramnote{通用反馈环。实线是前向通路：参考值与测量值在求和点相减得到误差，控制器按误差计算控制量，经执行器作用于被控对象；虚线是反馈通路：传感器测量输出并送回求和点。扰动从前向通路中途进入——这是反馈存在的理由：环路无法阻止扰动到来，只能在扰动造成偏差后把它纠正回去。每个框下方给出本章各环对应的实物。}

把四个案例按这套词汇对照一遍：

\begin{longtable}{L{0.13\linewidth}L{0.40\linewidth}L{0.40\linewidth}}
\toprule
词汇 & 机制 & 机器里的实例 \\
\midrule
被控对象 & 被维持状态的物理系统与被控量 & VRM 的 LC 输出级（电压）、VCO（相位）、磁头臂（位置）、硅片与散热器（温度） \\
传感器 & 把被控量变成可比较的信号 & 分压电阻网络、鉴相器、热敏二极管、伺服 burst 解调 \\
执行器 & 控制器推得动的物理入口 & buck 占空比、VCO 控制电压、音圈电机电流、风扇 PWM \\
参考值 & 目标值 & 带隙基准电压、参考时钟沿、温度阈值寄存器 \\
误差 & $e = r - \hat{y}$，控制器的唯一输入 & 误差放大器输入差、鉴相器脉宽、PES 数值 \\
前馈 / 反馈 & 测扰动提前补偿 / 测结果事后纠正 & VRM 负载线补偿；HDD 的 RRO 前馈表 + PID 闭环 \\
\bottomrule
\end{longtable}

这套词汇的价值在于可迁移。任何一个陌生的"环"，先问五个问题：被控量是什么、谁测量它、参考从哪里来、误差以什么物理形式存在、执行器推的是什么。五个问题答完，这个环的结构就清楚了，剩下的只是参数。

\topic{5.2\quad 稳定性直觉：纠正为什么会变成助推}

负反馈在直觉上是自我纠正的：偏差出现，环路施加反向作用，偏差缩小。但任何读过示波器上电源振荡波形的人都知道，负反馈环会振荡。问题出在\emph{时间}上：纠正信号不是瞬时生效的，它在环路里每经过一个储能环节就被延迟一点。

把纠正信号在环路里绕行一周的过程拆开看。每个储能元件——电感、电容、机械惯量、热容——都引入两种效应：信号幅度被衰减，相位被推迟。在最简模型里，控制理论把每个储能环节称为一个\hwkey{极点}（LC 谐振这样的一对储能元件实为一对共轭极点），一个极点最多贡献 $-90^{\circ}$ 的相位滞后。一个极点的环最多滞后 $90^{\circ}$，纠正方向只是"打了折扣"，永远不会反；两个以上极点串联时，总滞后可以向 $180^{\circ}$ 逼近。滞后达到 $180^{\circ}$ 意味着：输出偏高时，绕环一周回来的纠正信号不是把它压低，而是恰好同相地把它推得更高——负反馈在这个频率上变成了正反馈。

此时振荡是否发生，取决于\hwkey{环路增益}：扰动绕环一周，是被放大还是被缩小。若总滞后接近 $180^{\circ}$ 的那个频率上环路增益仍大于 1，任何微小的扰动绕环一周后都变得更大，再绕一周更大——振幅持续增长，直到撞上某个物理限制（电源轨、占空比上限）变成等幅振荡。若在那个频率上增益已经小于 1，扰动每绕一周都在衰减，环是稳定的。用推秋千类比最直观：在秋千荡到最高点时推一把，能量越来越大——这是滞后凑足半个周期、推变成助推；在秋千往回荡时推，越推越小——这才是纠正。相位滞后在低频时微不足道，纠正还是纠正；频率升高后滞后累积，同一套纠正逻辑就变成了助推。

工程上把这套直觉量化成两个裕量。增益降到 1 的那个频率称为\term{穿越频率}{crossover frequency}，它近似等于环路的响应速度——带宽。\term{相位裕度}{phase margin}是穿越频率处总滞后与 $180^{\circ}$ 的距离：裕度 $45^{\circ}$--$60^{\circ}$ 是教科书式的舒适区，小于 $30^{\circ}$ 的环阶跃响应有明显振铃，趋近 $0^{\circ}$ 就是趋近持续振荡。对称地，增益裕度是滞后达到 $180^{\circ}$ 处增益与 1 的距离，习惯要求 6 dB 以上。设计一个补偿网络，本质上就是安排环路增益的下降曲线：让它在相位滞后累积到危险区之前，先降得足够低。

积分环节在这里扮演双重角色。积分器——输出是输入的累积——天然贡献恒定的 $-90^{\circ}$ 滞后，且在低频段增益很高。环里放一个积分器，稳态误差会被积成零：只要误差不为零，积分就持续累积控制量，直到误差消失为止。这是它诱人的一面；代价是那固定的 $-90^{\circ}$ 把总滞后向 $180^{\circ}$ 推近了一大步。一阶惯性环节加积分环节的组合，是机器里最常见的环路骨架：波特图上，积分器先给出每十倍频 20 dB 的下降斜率，过了惯性环节的极点频率后两段叠加，斜率加大为每十倍频 40 dB——先缓后陡；相位曲线从 $-90^{\circ}$ 起步向 $-180^{\circ}$ 滑落，补偿网络要做的就是在相位曲线滑到危险区之前，把增益曲线按到 1 以下。本节不展开波特图的作法，记住趋势就够了：\hwkey{积分换精度，带宽换速度，裕度换平稳，三者互相抢}。

\begin{classicnote}
负反馈放大器是贝尔实验室的 Harold Black 在 1927 年为解决长途电话中继放大器的失真问题发明的；Harry Nyquist 随后在 1932 年给出了稳定性的严格判据。这套理论最初的"机器"就是模拟电子线路，后来被搬到机械、化工、生物系统上，1948 年 Wiener 的《控制论》把它确立为横跨工程的通用学科。计算机里的电压环、锁相环、伺服环，直接继承的是这条电话工程的血统——VRM 误差放大器的补偿网络画法，与九十年前的中继放大器是同一张图。
\end{classicnote}

稳定性取舍在时域里最直观的表现，是阶跃响应的三种形态。给环一个突然的目标变化（或一次负载突变），输出的回应方式暴露环路的全部性格：

\begin{pdfdiagram}[x=1cm,y=1cm]
  % 坐标轴
  \draw[BookRule, line width=0.5pt, ->] (0.5,0.5) -- (10.3,0.5);
  \node[hw label, anchor=north] at (10.3,0.38) {时间};
  \draw[BookRule, line width=0.5pt, ->] (0.5,0.5) -- (0.5,4.7);
  \node[hw label, anchor=south east] at (0.5,4.7) {输出};
  % 目标值虚线
  \draw[BookMuted, dashed, line width=0.5pt] (0.5,3.4) -- (10.1,3.4);
  \node[hw label, font=\tiny, anchor=east] at (10.05,3.62) {参考值 $r$};
  % 欠阻尼：超调后振荡收敛（BookRed）
  \draw[BookRed, line width=0.85pt] (0.5,0.5) .. controls (1.2,0.6) and (1.7,3.9) .. (2.3,4.1)
    .. controls (2.9,4.3) and (3.3,2.8) .. (4.0,2.78)
    .. controls (4.7,2.76) and (5.0,3.58) .. (5.7,3.6)
    .. controls (6.4,3.62) and (6.7,3.24) .. (7.4,3.26)
    .. controls (8.1,3.28) and (8.8,3.4) .. (10.1,3.4);
  \node[BookRed, font=\tiny, anchor=west] at (4.4,4.28) {欠阻尼：超调与振铃};
  % 临界阻尼：最快无超调（BookAccent）
  \draw[BookAccent, line width=0.85pt] (0.5,0.5) .. controls (1.6,0.6) and (2.5,3.1) .. (3.7,3.34)
    .. controls (5.0,3.42) .. (7.2,3.4) -- (10.1,3.4);
  \node[BookAccent, font=\tiny, anchor=west] at (7.9,2.6) {临界阻尼：最快且无超调};
  % 过阻尼：缓慢逼近（BookTeal）
  \draw[BookTeal, line width=0.85pt] (0.5,0.5) .. controls (2.5,0.55) and (4.0,1.5) .. (5.5,2.4)
    .. controls (7.0,3.0) and (8.5,3.3) .. (10.1,3.38);
  \node[BookTeal, font=\tiny, anchor=west] at (5.6,1.85) {过阻尼：平稳但迟缓};
\end{pdfdiagram}
\diagramnote{同一阶跃目标下三种阻尼形态的响应。欠阻尼环（相位裕度小）上升最快，但冲过目标后振铃数次才停稳；过阻尼环（裕度大、带宽低）不越雷池一步，却慢得多；临界阻尼是不超调前提下最快的形态。工程上常故意留一点超调（相位裕度 $45^{\circ}$ 上下，超调约 15\%--25\%），用少量振铃换取明显更快的响应。}

三种形态对应三种设计失败或妥协。欠阻尼：补偿网络给的相位裕度太小，环在追赶目标时刹不住车。过阻尼：积分增益或带宽取得太保守，环对任何变化都慢半拍——负载突变时下陷又深又久。临界阻尼是理论最优点，实际设计宁可靠近它而偏一点欠阻尼，因为元件公差和温度会把参数推来推去，贴着临界值设计的环在批量产品中总有一部分滑进深欠阻尼区。

\begin{warningbox}
"接成负反馈就稳定"是最常见的误解。反馈的极性只在低频保证纠正方向正确；频率升高后各储能环节的相位滞后累积，负反馈会在某个频率上变成正反馈。稳定与否取决于环路增益和相位在整个频率轴上的配合，不取决于接线端的正负号。第二个常见误解是"积分增益越大、稳态越准"：积分确实把稳态误差积成零，但它的 $-90^{\circ}$ 滞后同时压缩相位裕度，积分调过头，环会以低频大幅振荡的方式"记住"这个错误。
\end{warningbox}

\topic{5.3\quad 案例一：VRM 电压环——二阶对象与比例-积分补偿}

本章第四节讲过 VRM 的结构：多相 buck 把 12 V 降到核心所需的 1 V 上下，误差放大器比较输出分压与参考，调整占空比。现在用控制理论重读这个环。

被控对象是二阶的。buck 输出级的电感与输出电容各是一个储能元件，LC 组合存在一个谐振频率（典型几十 kHz）：低于谐振频率，输出电压忠实地跟随占空比；接近谐振频率，相位急剧滑落、增益出现尖峰——若不加处理，环在谐振频率附近凑齐振荡条件。误差放大器外围的 RC 网络就是为这件事存在的：它实现\hwkey{比例-积分补偿}，积分支路在低频提供高增益，把直流输出误差积成零（这就是稳压精度的来源）；比例支路与精心放置的零点在穿越频率附近托住相位，把相位裕度维持在 $45^{\circ}$--$60^{\circ}$。补偿网络的几个电阻电容取值，就是这个环的相位裕度的物理形态。

比例-积分之外还可以加上微分环节：微分正比于误差的变化率，误差还在加速恶化时它就提前加大纠正量，相当于给环路装上预判；在频域上它提供超前相位，能部分抵消储能环节带来的滞后、抑制振荡——这正是 4.2 节"微分环节抑制振荡"的含义。只是微分对高频噪声同样敏感，会把输出纹波放大进控制量，实际 VRM 多以 PI 为主、微分为辅。

负载瞬态是观察这个环的最佳实验。CPU 从空闲跳进全负载，电流边沿的 di/dt 可达几十 A/ns，一次完整的负载阶跃在几十纳秒到微秒量级内完成（电源时钟章第一节算过这笔账）。环路带宽折算成响应时间是 $\mu$s 级，最初的下跌环路根本来不及参与：电压先按输出电容的等效串联电阻跌一个台阶，再因电容放电继续下滑，直到环路把占空比提上来才止步回升。下陷的深度近似为 $\Delta V \approx \Delta I \times Z_{out}$——$Z_{out}$ 是输出端看进去的等效阻抗，由电容参数与环路增益共同决定；恢复的快慢由环路带宽决定，穿越频率越高，环越早察觉、越早纠正。电源时钟章第一节"去耦电容按时间尺度分层"说的正是同一个事实的另一面：ns 级靠片上电容、$\mu$s 级靠板级电容，再慢才轮到环路——每一层只负责自己来得及管的时间尺度。

\begin{pdfdiagram}[x=1cm,y=1cm]
  % 负载电流阶跃
  \draw[BookAccent, line width=0.85pt] (0.5,5.0) -- (4.0,5.0) -- (4.0,6.2) -- (10.3,6.2);
  \node[hw label, anchor=east] at (0.4,5.6) {负载电流};
  \draw[BookAmber, line width=0.5pt, <->] (3.7,5.0) -- (3.7,6.2);
  \node[hw label, font=\tiny, anchor=east] at (3.55,5.6) {$\Delta I$};
  % 输出电压：基线 + 下陷 + 恢复
  \draw[BookMuted, dashed, line width=0.5pt] (4.0,3.0) -- (10.3,3.0);
  \draw[BookTeal, line width=0.85pt] (0.5,3.0) -- (4.0,3.0) -- (4.15,2.55)
    .. controls (4.4,1.95) and (4.8,1.8) .. (5.3,1.95)
    .. controls (5.8,2.1) and (6.2,3.05) .. (6.7,3.08)
    .. controls (7.2,3.11) and (7.6,3.0) .. (8.3,3.0) -- (10.3,3.0);
  \node[hw label, anchor=east] at (0.4,3.0) {输出电压};
  % 下陷深度标注
  \draw[BookAmber, line width=0.5pt, <->] (5.45,3.0) -- (5.45,1.95);
  \node[hw label, font=\tiny, anchor=west] at (6.9,2.45) {下陷 $\approx \Delta I \times Z_{out}$};
  % 恢复时间标注
  \draw[BookMuted, line width=0.5pt, <->] (4.0,1.3) -- (8.3,1.3);
  \node[hw label, font=\tiny, anchor=north] at (6.15,1.18) {恢复时间：由环路带宽决定};
  % 时间轴
  \draw[BookRule, line width=0.5pt, ->] (0.5,0.7) -- (10.3,0.7);
  \node[hw label, anchor=north] at (10.3,0.58) {时间};
\end{pdfdiagram}
\diagramnote{负载电流阶跃下的 VRM 输出电压。最初的陡跌由输出电容的寄生电阻决定，环路来不及参与；随后的下滑与回升才是环路带宽的表演时间。带宽越高的环下陷越浅、恢复越快，但带宽受 LC 谐振与相位裕度限制，不能无限推高。}

多相并联是执行器层面的优化。第四节讲过 $N$ 相 buck 错相 $360^{\circ}/N$ 驱动：各相纹波电流峰谷互错，在输出电容上叠加时大部分抵消，等效纹波频率变成 $N$ 倍、幅度大减。用控制理论的话说，这是把一个大执行器拆成 $N$ 个错峰工作的小执行器：各相电感器相同，但 $N$ 相并联后等效电感为 $L/N$，相电流对占空比的响应更快，环路可以安全地取更高的带宽；输出电容的负担也随之减轻。软件碰到这个环的方式是 4.2 节讲过的 VID 通路：CPU 把目标电压编码发给 VRM 控制器，控制器平移自己的参考值，环路自动把输出调过去。DVFS 里的"调压"没有一根模拟控制线，软件改的就是这个参考值——完整的软硬协同故事见 5.5。

\topic{5.4\quad 案例二：PLL 与 DLL——锁相环为什么是 II 型系统}

电源时钟章第二节画过 PLL 的五环结构，存储器件结构章第三节讲过 PHY 里的 DLL。用本节词汇重读：鉴相器是传感器，输出正比于相位差的脉冲；电荷泵加环路滤波器是控制器，把误差脉冲积成平缓的控制电压；VCO 是执行器，控制电压改变振荡频率；被控量是输出时钟的相位。

这个环有一个在控制理论里很特殊的性质：它天然带两个积分器。第一个是 VCO 本身——控制电压决定的是\emph{频率}，而频率是相位的导数，相位是频率对时间的积分。被控量是相位、执行器输出的是频率，这中间隔着一次免费的积分。第二个积分器在环路滤波器里：理想电荷泵把相位误差脉冲积分成控制电压。环路里串着两个积分器的系统称为 \hwkey{II 型系统}，它有两个直接后果。其一，锁定后静态相位误差为零：只要参考沿与反馈沿之间还有恒定的相位差，第二个积分器就持续累积控制电压，直到相位差消失。其二，它能无静差地跟踪频率阶跃：参考频率跳变后，环路重新锁定且相位仍然对齐。这就是为什么 PLL 是时钟倍频的标准方案——它对"准确"的承诺是结构保证的，不靠手工调参。

DLL 少一个积分器。它的执行器是数控延迟线：控制量直接就是相位（延迟），不需要经过"频率积分成相位"那一道。环路里只剩一个积分器，是 I 型系统。少一个积分器意味着少 $-90^{\circ}$ 的固有滞后，相位裕度天然宽裕，因此 DLL 锁定快、不易振铃——这正是存储器件结构章第三节说"DLL 锁定快、引入的抖动小"的控制论原因。代价同样来自结构：延迟线不产生新频率，DLL 只能搬移相位，不能倍频。

抖动在控制理论里是环路的\hwkey{噪声传递}问题。参考时钟上的噪声进入环路后按低通特性传到输出：带宽以内的慢起伏，环路跟得上、照单全收；带宽以外的快抖动，环路来不及反应。VCO 自身的噪声正好相反，按高通特性出现：环路能纠正慢漂移，管不了快起伏。环路带宽因此是两类噪声的分界旋钮——带宽越低，参考噪声滤得越干净，但 VCO 噪声露出来的频段越宽。电源时钟章第二节讲"用 PLL 给抖动的参考洗干净"时把带宽设低，选的就是这个权衡的一端。软件与这个环的唯一接触点是锁定状态：启动代码先配好分频比，再轮询 locked 标志，确认环路收敛之后才切换时钟源——一个 II 型系统的收敛时间，就这样站在每台机器开机流程的关键路径上。

\topic{5.5\quad 案例三：DVFS 与热管理——软硬协同的慢环}

前两个案例的环路带宽以 kHz--MHz 计，这一节看频谱的另一端：以秒和分钟计的热环。被控对象是硅片、导热界面材料与散热器组成的热系统：硅片功耗经热阻流向散热器，各环节的热容把温度变化拖慢——芯片结温的时间常数以秒计，散热器整体以分钟计。被控对象慢，环也快不了，而且不必快：一个热时间常数 10 秒的系统，用 1 ms 的环路去管纯属浪费。这是控制设计里一条朴素的原则：\hwkey{环路的响应速度只需匹配被控对象与扰动的速度}。

环的组成：传感器是分布在芯片热点附近的数字温度传感器（DTS）和电流监测电路；执行器有三个层次——频率（PLL 重锁或时钟分频）、电压（经 SVID 改 VRM 参考值，即 5.3 的 VID 通路）、以及功耗门控（整核时钟门控）；控制器的位置则经历了一次意味深长的迁移。早期方案（Intel SpeedStep 时代）由操作系统当控制器：OS 每隔 10--100 ms 评估一次负载，从 ACPI 定义的 P-state 表里挑一档，写 MSR 请求切换。这个环的控制周期是几十毫秒级，对编译、视频解码这类持续负载够用，但对"用户点开一个网页，渲染线程突然忙 20 ms"这种脉冲式负载完全跟不上——等 OS 察觉时活已经干完了，频率白提一次；或者更糟，升频姗姗来迟，用户先感到卡顿。Intel 从 Skylake 起把控制器搬进硬件（Speed Shift，硬件管理 P-state）：硬件以约 1 ms 的周期自主评估负载并调频，OS 退居二线，只写下一个性能偏好提示（EPP），从"每步都指挥"变成"只定基调"。控制周期从几十毫秒缩到一毫秒，环能跟踪的负载变化速度提高了一个数量级以上——这就是控制周期决定跟踪能力的直接演示。

风扇控制是同一个环在机箱尺度的延伸，结构更简单：一条\hwkey{风扇曲线}，本质是分段的\hwkey{比例控制}——温度低于阈值给固定低速，中间区间转速随温度线性上升（斜率就是比例增益），超过上限全速。纯比例控制会有稳态误差，风扇环不在乎：目标不是把温度钉在某一度，而是"别过热也别太吵"，误差本身被折算进了曲线设计。曲线上还有一个控制理论的经典元件：\term{迟滞}{hysteresis}。升温和降温走不同的阈值——比如升到 60$^{\circ}$C 才提速，降到 55$^{\circ}$C 才降速——否则温度在阈值附近抖动时，风扇会每秒切换一次转速，噪声反而更恼人。迟滞是用状态记忆换取决策稳定性，它是所有"阈值切换"类控制的必备零件。

\begin{lstlisting}
# 硬件监控芯片（hwmon）暴露给软件的接口：温度与风扇转速都是只读文件
$ cat /sys/class/hwmon/hwmon2/temp1_input     # 结温，单位 0.001 摄氏度
62000
$ cat /sys/class/hwmon/hwmon2/fan1_input      # 风扇实测转速 RPM
1450

; EC 固件里的风扇环（约每 100 ms 执行一次）
fan_loop:
  T = read_dts()
  if T < 50 then
      pwm = 30%                       ; 低速段（固定）
  else if T < 80 then
      pwm = 30% + (T - 50) * 2.3%     ; 比例段，斜率=增益
  else
      pwm = 100%                      ; 全速段（饱和）
  pwm = hysteresis(pwm)               ; 升降温走不同阈值
  write_pwm(pwm)
\end{lstlisting}

\begin{pdfdiagram}[x=1cm,y=1cm]
  % 坐标轴
  \draw[BookRule, line width=0.5pt, ->] (0.8,0.8) -- (10.4,0.8);
  \node[hw label, anchor=north] at (10.4,0.68) {温度};
  \draw[BookRule, line width=0.5pt, ->] (0.8,0.8) -- (0.8,4.4);
  \node[hw label, anchor=south east] at (0.8,4.4) {风扇 PWM};
  % 参考刻度
  \draw[BookMuted, dashed, line width=0.4pt] (4.0,0.8) -- (4.0,1.6);
  \node[hw label, font=\tiny, anchor=north] at (4.0,0.7) {50$^{\circ}$C};
  \draw[BookMuted, dashed, line width=0.4pt] (7.5,0.8) -- (7.5,3.9);
  \node[hw label, font=\tiny, anchor=north] at (7.5,0.7) {80$^{\circ}$C};
  \draw[BookMuted, dashed, line width=0.4pt] (0.8,3.9) -- (7.5,3.9);
  \node[hw label, font=\tiny, anchor=east] at (0.7,3.9) {100\%};
  % 风扇曲线：低速恒定 -> 比例段 -> 全速（实线为升温路径，虚线为降温路径）
  \draw[BookAccent, line width=0.9pt] (0.8,1.6) -- (4.0,1.6) -- (7.5,3.9) -- (10.2,3.9);
  \draw[BookAccent, line width=0.7pt, dashed] (0.8,1.6) -- (3.5,1.6) -- (7.0,3.9) -- (10.2,3.9);
  \draw[BookAccent, line width=0.5pt, <->] (6.0,2.91) -- (6.0,3.24);
  \node[BookAccent, font=\tiny, anchor=west] at (6.15,3.08) {迟滞};
  \node[hw label, font=\tiny, text width=3.4cm] at (9.0,2.35) {实线：升温路径；虚线：降温路径，两路径的间距即迟滞区间};
  \node[BookAccent, font=\tiny] at (2.3,1.85) {低速恒定};
  \node[BookAccent, font=\tiny, anchor=west] at (5.1,1.92) {比例段：斜率即增益};
  \node[BookAccent, font=\tiny, anchor=south] at (8.8,3.98) {全速（饱和）};
\end{pdfdiagram}
\diagramnote{风扇曲线：分段比例控制。低温段固定转速避免频繁启停；中段转速随温度线性上升，斜率就是比例增益，斜率越大控温越紧、噪声越敏感；高温段进入饱和，全速运行。50$^{\circ}$C 与 80$^{\circ}$C 两个转折点附近各有迟滞区间，防止温度抖动引起转速来回切换。}

软件在三条路径上碰到这个慢环。其一，读接口：hwmon 子系统把温度、风扇转速导出为 sysfs 文件，\code{sensors} 命令读的就是它们。其二，写接口：\code{cpufreq} 子系统的 governor（performance/powersave/schedutil）决定 OS 侧的策略，Speed Shift 平台上策略退化为一个 EPP 提示值。其三，越界接口：温度越过临界点时，硬件不等任何软件——PROCHOT 信号直接拉低频率，再不行就 THERMTRIP 强制断电。安全相关的最后一级永远做在硬件里，因为热失控时不能假设软件还活着。

\topic{5.6\quad 案例四：内存训练——没有微分方程时的控制}

前三个案例的环有一个共同前提：被控对象可以写成连续的、近似可微的模型——电压对占空比、相位对控制电压、温度对功耗，都是光滑变化的量，误差信号连续，纠正量也连续。内存训练面对的世界完全不同。存储器件结构章第三节讲过训练的细节：PHY 扫描每根数据线的采样延迟与判决阈值，DRAM 回报的只有 PASS/FAIL 一个比特。这里没有"误差有多大"，只有"过与没过"；眼图的形状由走线、过孔、封装、温度共同决定，无法写成一个可以求导的传递函数。连续控制的所有工具——增益、相位、补偿——在这里都无从下手。

工程上的替代方案是把"调节"换成\hwkey{系统辨识}加\hwkey{参数搜索}。第一步，测量对象：控制器把延迟和 VREF 扫成二维网格，每个格点上做一次读写测试，PASS/FAIL 填满整张表——这张表就是被控对象的特性画像，它在时间轴上的投影就是眼宽，在电压轴上的投影就是眼高。第二步，在画像上做优化：找 PASS 区域的中心点，把采样参数定在那里——中心是各方向裕量最大的位置，对日后的温度漂移和老化最稳健。搜索策略本身可以讲究：粗扫定位大区、细扫逼近边界、二分查找精确边沿，都是标准优化技术在这个离散问题上的应用。这与连续环的区别是根本性的：连续环\emph{持续地}用误差修正输出，训练是\emph{一次性地}把参数标定好，然后开环运行——它更像秤的校准，而不是定速巡航。

但这条边界正在移动。标定结果会随温度漂移：硅片的延迟随温度变化，标定时找到的眼中心，在机器跑热之后可能偏离几十 ps。LPDDR 平台与部分 DDR5 平台的应对是把一次性校准升级为慢速的持续调节：运行中周期性重训练；JEDEC DDR5 规定的则是周期性 ZQ 校准，按温度变化重新标定驱动阻抗（存储器件结构章第三节）。训练的范式由此从"出厂一次"走向"运行中微调"，离散搜索与连续反馈在这个时间尺度上汇合了。同一趋势也出现在高速串行链路里：PCIe/USB 接收端的连续时间均衡器是模拟反馈环，判决反馈均衡器（DFE）则每个比特做一次离散调整——一条链路上两种范式并存，各管一段频率。

\topic{5.7\quad 两种控制范式总结}

把本章和此前各章的"控制"清点一遍，可以分成两大族。一族是\hwkey{连续反馈环}：被控量是连续的物理量，误差信号连续存在，纠正量连续可调，分析工具是环路增益、相位与阶跃响应——VRM 电压环、PLL、HDD 磁头伺服（其 PID 细节见存储器件结构章第六节，此处不再重复）、风扇与热环都在此列。另一族是\hwkey{离散状态机}：被控对象是协议与资源的状态，输入是事件与应答，纠正手段是状态转移，分析工具是状态图与时序图——内存控制器的 bank 状态机（存储器件结构章第四节）、DMA 通道的传输状态机（主存与总线章）、PCIe 枚举与链路训练、内存训练，都是这一族。

\begin{longtable}{L{0.16\linewidth}L{0.38\linewidth}L{0.38\linewidth}}
\toprule
维度 & 连续反馈环 & 离散状态机 \\
\midrule
被控对象 & 连续物理量：电压、相位、位置、温度 & 协议与资源状态：bank 行开闭、传输进度、链路状态 \\
反馈信号 & 连续误差：电压差、相位差、位置误差 & 离散事件与应答：PASS/FAIL、完成中断、状态位 \\
调节手段 & 连续纠正量：占空比、控制电压、电流 & 状态转移：发命令、改配置、切换模式 \\
时间尺度 & 由对象时间常数决定，ns（VRM/PLL）到分钟（热） & 由协议节拍决定，一个总线周期到一次训练 \\
典型失效 & 失稳：振荡、超调、锁定丢失 & 逻辑错误：死锁、超时、挂死在某个状态 \\
分析工具 & 环路增益、相位裕度、阶跃响应 & 状态图、时序图、超时与重试机制 \\
本章实例 & VRM 电压环、PLL/DLL、DVFS 热环、风扇 & 内存训练、bank 状态机、电源时序、链路枚举 \\
\bottomrule
\end{longtable}

两族的分界不在硬件与软件之间，而在\emph{被控对象的数学性质}之间：能写出连续模型的用环，写不出的用状态机。真实的控制器几乎总是两者的复合体——VRM 芯片里，电压环是连续反馈，上电时序、过流保护、多相切相是状态机；内存控制器里，调度是状态机，DLL 与周期性重训练是环；HDD 里，寻道是状态机切换三段轮廓，磁道跟踪是连续伺服。整机之所以可靠，不是因为哪一族更优越，而是每个问题都被交给了与它数学性质匹配的那一族。

读懂一台机器里所有的控制环之后，"CPU 如何控制整台机器"这个问题就有了第二层答案：CPU 并不直接控制机器，它配置的是一大批自治的控制环与状态机——设定参考值、读取状态字、接收中断——真正的控制在环里，以各自的物理节奏持续进行，无论 CPU 此刻在不在看它们。
\section{六、软件侧的控制：驱动、中断与调度}

本章到目前为止讨论的控制环都运行在硅片和固件里：VRM 在纳秒尺度上稳住电压，PLL 在皮秒尺度上对齐相位，内存控制器在每个时钟周期决定发哪条命令。这些环对软件完全透明——操作系统看不到它们，也不该看到。但整台机器还有另一半控制问题：设备谁来找、找到了谁来管、设备有事了谁被叫醒、CPU 时间按什么规则分。这一层的控制周期在微秒到秒之间，控制器就是操作系统本身。本节回答三个问题：操作系统怎么发现设备并把驱动的代码接到设备寄存器上，设备事件怎么沿中断链一路上行到某个等待的进程，以及定时器中断怎样把调度器变成一个离散时间控制器。

\begin{keyidea}
软件侧的控制和硅片里的控制环是同一套结构，只是时间尺度不同。设备寄存器是执行器的接口，中断是传感器送到 CPU 的信号，驱动是环路的控制律实现。把它们当作控制环来读，驱动代码里"读状态—发命令—等完成"的固定套路就不是经验习惯，而是闭环的必要环节。
\end{keyidea}

\topic{6.1\quad 操作系统怎么发现设备}

驱动要能控制设备，前提是操作系统先知道设备存在、知道它的寄存器在地址空间的哪个窗口。发现设备有三条途径，分别对应三种硬件现实。

第一条途径是 \term{ACPI 表}{Advanced Configuration and Power Interface}。PC 和服务器上，固件在内存里放一组表，其中 DSDT 等表用 ACPI 机器语言描述主板上"焊死"的设备——嵌入式控制器、温度传感器、电源按钮——以及每个设备占用的 I/O 端口、MMIO 窗口和中断线。操作系统启动时解析这套命名空间，为每个设备节点建立内部对象，再按节点里的硬件标识符（HID）去驱动库里找匹配的驱动。第二条途径是\term{设备树}{device tree}。嵌入式 ARM、RISC-V 平台没有固件代为描述的传统，改为 bootloader 把一份二进制设备树（DTB）随内核一起传入内存，里面逐节点列出 SoC 内部总线和挂在上面的控制器；驱动按节点里的 compatible 字符串声明"我能管这类设备"，内核按字符串配对。第三条途径是 \term{PCI 枚举}{PCI enumeration}：PCI/PCIe 设备是可插拔的，无法事先写进任何静态表，但协议本身提供了发现手段——每个设备的开机配置空间里都有厂商号、设备号，主机按总线号逐设备读出这些编号，再读出 BAR 寄存器声明的 MMIO 窗口大小，把窗口分配进物理地址空间。驱动携带一张自己支持的编号表，枚举到一个设备就比对一次，命中即接管。经典芯片与平台章讲过平台启动时枚举在开机流程中的位置，这里关心的是三条途径的共同产物：一个"设备对象 + 一组已分配的地址窗口 + 一段被调用的驱动入口函数"。

\begin{longtable}{L{0.12\linewidth}L{0.32\linewidth}L{0.20\linewidth}L{0.27\linewidth}}
\toprule
途径 & 信息载体 & 典型平台 & 驱动匹配的入口 \\
\midrule
ACPI 表 & 固件放置在内存中的表，描述设备存在与资源分配 & PC、服务器 & 解析命名空间，按硬件标识符（HID）找驱动 \\
设备树 & bootloader 随内核传入的二进制描述（DTB） & 嵌入式 ARM、RISC-V & 按 compatible 字符串配对驱动 \\
PCI 枚举 & 总线协议自身：配置空间里的厂商号、设备号与 BAR & 所有含 PCI/PCIe 的平台 & 按驱动自带的编号表逐项比对 \\
\bottomrule
\end{longtable}

三条途径的分工由硬件的可插拔性决定：焊在主板上的设备写进静态描述（ACPI 或设备树），插槽上的设备靠协议自己报告（PCI 枚举）。匹配成功之后，驱动拿到的是设备窗口的物理地址，还不能直接用——x86-64 内核运行在虚拟地址上，设备窗口必须先映射进来。\code{ioremap} 完成这件事：它在内核虚拟地址空间中建立一段到该物理窗口的映射，并给这段映射打上设备内存属性（不可缓存、不可合并、写顺序受保护），这些属性正是主存与总线章反复强调的 MMIO 访问约定。

映射建立之后，驱动读写设备寄存器和读写普通内存的代码在指令层面完全一样——都是 \code{mov}。下面是一个最小设备驱动的典型骨架，左侧是 C 形式，右侧是它编译后对应的 x86-64 指令形态：

\begin{lstlisting}
/* base 是 ioremap 返回的内核虚拟地址；
 * SR: 状态寄存器 (bit0 = BUSY)，CR: 命令寄存器，IE: 中断使能 */
while (readl(base + SR) & BUSY)      /* 读 SR 判忙 */
    cpu_relax();
writel(CMD_START, base + CR);        /* 写 CR 发命令 */
writel(IRQ_DONE,  base + IE);        /* 开中断，然后睡眠等待 */

/* 上面三行编译后就是普通的 mov，与读写内存同一类指令：
 * mov  eax, [rbx]          ; 读 SR
 * test eax, 1
 * jnz  .busy
 * mov  dword [rbx+4], 1    ; 写 CR
 * mov  dword [rbx+8], 1    ; 写 IE                          */
\end{lstlisting}

代码和内存访问长得一样，语义却完全不同：读状态寄存器拿到的是设备此刻的硬件事实，写命令寄存器是推了一下设备内部的状态机。主存与总线章已经指出关键一点——MMIO 写返回只说明命令到达设备接口，不代表动作完成。完成的通知有两条路：驱动死等状态位（轮询），或者像上面这段代码一样开好中断去睡眠，把 CPU 让给别的进程。慢设备上轮询是纯粹的浪费，于是需要中断把"完成"这件事主动送回来——这是下一小节的主题。

\topic{6.2\quad 中断处理链：从设备事件到被唤醒的进程}

主存与总线章从硬件角度讲过中断链路：事件在设备内锁存成状态位并拉起请求线，中断控制器按屏蔽位和优先级裁决，CPU 在指令边界采样、保存现场、查表跳转。这里把同一条链按现代 x86-64 平台补全，并讲清链路的软件后半段。

现代平台上有两种中断投递方式。老方式是引脚：设备拉起一条物理中断线，进入 \term{IO/APIC}{I/O Advanced Programmable Interrupt Controller}——现代平台上它集成在 PCH 内（早期主板上是独立的 82093AA 芯片），\hwevidence{典型实现有 24 条重定向输入}，每条对应一张重定向表项，写着这条线应翻译成几号向量、投给哪个核。新方式是 \term{MSI}{message signaled interrupt}：设备不拉线，而是向一段特殊地址窗口做一次内存写事务，写的内容里直接带着向量号和目标核。MSI 省掉了引脚和共享线的竞争，还能让一个设备拥有多个向量——每核一个队列一个向量，外设队列章里 NVMe 的每核发中断靠的就是它。两条路殊途同归，终点都是目标核的\term{本地 APIC}{Local APIC}：它按优先级决定何时把向量递交给 CPU。

CPU 在指令边界接收向量后做的事由硬件完成：保存返回地址和标志，按向量号查\term{中断描述符表}{interrupt descriptor table, IDT}——一张 256 项的表，每项是一个处理函数的入口——跳到对应表项。从这里开始是软件。处理函数即 \term{ISR}{interrupt service routine, 中断服务例程}，它最急的工作只有两件：按设备规则清掉中断源（否则同一事件会反复触发），以及把设备里等着的少量数据取到内存安全的地方。做完这两件，ISR 就应该退场，把剩下的慢活——协议解析、缓冲区整理、唤醒逻辑——推迟给下半部：\term{软中断}{softirq}、\code{tasklet} 或者内核线程，由它们在允许睡眠、允许被抢占的上下文里慢慢做完，最后唤醒睡眠在等待队列上的进程。整条链如下：

\begin{pdfdiagram}[x=1cm,y=1cm]
  % --- 主链：设备 → IO/APIC → Local APIC → CPU → ISR ---
  \node[hw compute, minimum width=1.7cm] (dev) at (0,0) {设备\\事件/状态位};
  \node[hw control, minimum width=1.7cm] (ioapic) at (2.5,0) {IO/APIC\\重定向表};
  \node[hw control, minimum width=1.7cm] (lapic) at (5.0,0) {Local APIC\\IRR/ISR};
  \node[hw compute, minimum width=1.7cm] (cpu) at (7.5,0) {CPU 核\\查 IDT};
  \node[hw control, minimum width=1.7cm] (isr) at (10.0,0) {ISR 上半部\\清源/取数};
  % --- 下半部与进程（下方一行） ---
  \node[hw compute, minimum width=1.7cm] (bh) at (10.0,-2.2) {下半部\\softirq/线程};
  \node[hw state, minimum width=1.7cm] (proc) at (7.5,-2.2) {等待进程\\被唤醒};
  % --- 主链箭头 ---
  \draw[hw bus] (dev) -- node[above,hw label]{中断线} (ioapic);
  \draw[hw bus] (ioapic) -- node[above,hw label]{向量} (lapic);
  \draw[hw bus] (lapic) -- node[above,hw label]{投递} (cpu);
  \draw[hw bus] (cpu) -- node[above,hw label]{保存现场} (isr);
  \draw[hw bus] (isr) -- node[left,hw label,fill=white,inner sep=1pt]{推迟慢活} (bh);
  \draw[hw bus] (bh) -- node[below,hw label]{唤醒} (proc);
  % --- MSI 直通路径：设备直接写消息到 Local APIC ---
  \draw[hw bus] (dev.south) -- (0,-1.1) -- (5.0,-1.1) -- (lapic.south);
  \node[hw label, fill=white, inner sep=1pt] at (2.5,-1.1) {MSI：向 APIC 地址写消息};
  % --- EOI 返回路径 ---
  \draw[hw return] (isr.north) -- (10.0,1.1) -- (5.0,1.1) -- (lapic.north);
  \node[hw label, fill=white, inner sep=1pt] at (7.6,1.1) {EOI：恢复投递};
\end{pdfdiagram}
\diagramnote{中断处理链。实线是事件上行路径：引脚方式经 IO/APIC 重定向，MSI 方式由设备直接向 Local APIC 的地址窗口写消息；CPU 按向量查 IDT 进入 ISR，上半部只做清源和取数，慢活推迟给下半部，最后由下半部唤醒等待的进程。虚线是返回路径：ISR 向 Local APIC 发 EOI，控制器才恢复同级和更低优先级请求的投递。}

\begin{warningbox}
"中断处理函数里可以慢慢处理"是最常见的驱动错误。ISR 运行在中断上下文：它没有自己的进程身份，不能睡眠——一旦睡眠，没有任何东西能把它再调度回来；它还可能正压着同级和更低优先级的中断（EOI 之前这些请求都在等待）。在 ISR 里做耗时操作，后果不是这个设备变慢，而是整台机器的中断延迟被拉长，别的设备开始丢数据。规则只有一条：上半部做最急的事，其余全部推迟。
\end{warningbox}

\topic{6.3\quad 定时器中断驱动调度：操作系统的心跳}

中断链上最特殊的一个"设备"是定时器：它不报告外部事件，而是周期性地打断 CPU，给操作系统一个固定的决策节拍。现代 x86-64 平台上，这个节拍来自每个核本地的 \term{Local APIC 定时器}{Local APIC timer}——它按可编程的周期产生中断，内核把它配置成每秒 \hwtiming{250 或 1000 次}（内核配置项 HZ），即每 \hwtiming{4 ms 或 1 ms} 一个节拍，称为一次 \hwkey{tick}。

每个 tick 的处理路径很短，做的却是控制环里"传感器"和"决策"两件事。传感器一侧：更新系统计时，并给当前正在运行的进程累计这一段实际运行时间。决策一侧：检查当前进程的时间片是否用完、是否有更高优先级的进程变得可运行；需要换人就设置一个重调度标志，在中断返回用户态之前检查这个标志、进入调度器。调度器选中下一个进程，执行上下文切换——把当前进程的寄存器现场存进它的进程描述符，把下一个进程的现场装回 CPU。

把这个结构用控制语言写出来，它就是一台教科书式的离散时间控制器：被控量是各进程实际获得的 CPU 时间分配，传感器是定时器中断和运行时间累计，控制器是调度算法，执行器是上下文切换。与硅片里的连续环不同，它按固定节拍采样、按节拍决策——控制周期就是 tick。调度算法本身以 CFS（Completely Fair Scheduler，完全公平调度器）为例：它给每个进程维护一个虚拟运行时间 vruntime，把实际运行时间按优先级加权折算后累加，每次调度选 vruntime 最小的进程运行，使长期来看各进程的 CPU 份额收敛到公平值。至此，整机控制链的最后一环闭合了：连"CPU 这台执行器把时间分给谁"，也是一个有传感器、有控制律、有执行器的反馈环。

\begin{pdfdiagram}[x=1cm,y=1cm]
  \node[hw control, minimum width=1.75cm] (timer) at (0,0) {Local APIC\\定时器};
  \node[hw compute, minimum width=1.75cm] (tick) at (2.7,0) {tick 处理\\计时/累计};
  \node[hw control, minimum width=1.75cm] (sched) at (5.4,0) {调度器\\比较/选择};
  \node[hw compute, minimum width=1.75cm] (ctxsw) at (8.1,0) {上下文切换\\（执行器）};
  \node[hw state, minimum width=1.75cm] (proc) at (10.8,0) {当前进程\\（被控对象）};
  \draw[hw bus] (timer) -- node[above,hw label,font=\tiny]{节拍中断} (tick);
  \draw[hw bus] (tick) -- node[above,hw label,font=\tiny]{调度标志} (sched);
  \draw[hw bus] (sched) -- node[above,hw label,font=\tiny]{选下一个} (ctxsw);
  \draw[hw bus] (ctxsw) -- node[above,hw label,font=\tiny]{切换} (proc);
  \draw[hw return] (proc.south) -- (10.8,-1.2) -- (2.7,-1.2) -- (tick.south);
  \node[hw label, font=\tiny, fill=white, inner sep=2pt] at (6.75,-1.2) {实测运行时间：每个 tick 累计给当前进程};
\end{pdfdiagram}
\diagramnote{调度器作为离散时间控制器。Local APIC 定时器以固定节拍产生中断（采样）；tick 处理完成计时并把这一段运行时间累计给当前进程（反馈）；调度器比较后选中下一个进程；上下文切换把决策施加到 CPU 上。控制周期等于 tick 间隔，被控量是各进程实际获得的 CPU 时间。}

\topic{6.4\quad 边界原则：哪些环必须在硬件里，哪些可以放软件}

到这一步，本章出现过的控制环已经分布在从皮秒到秒的十二个数量级上。它们各自落在硬件、固件还是软件里，不是历史偶然，而是一条可计算的边界。判据有两条。第一条来自控制理论：\hwkey{控制周期必须比被控对象的时间常数快至少一个量级}，否则扰动已经发展完成，纠正动作才刚到达，环路不是稳住对象而是跟着振荡。第二条来自软件的物理下界：软件环最短的周期是一次中断加一次上下文切换，典型在微秒量级；再慢的策略性软件（用户态守护进程）响应在毫秒到秒。凡是被控对象的时间常数小于这个下界的环，没有任何软件选项，只能放在硅片里。

\begin{longtable}{L{0.22\linewidth}L{0.24\linewidth}L{0.20\linewidth}L{0.26\linewidth}}
\toprule
控制环 & 被控对象的时间常数 & 控制周期（量级） & 实现位置 \\
\midrule
VRM 电压环 & 负载瞬变，纳秒级 & 微秒级 & 硬件：VRM 控制器（模拟环） \\
PLL / DLL 相位环 & 相位漂移，皮秒至纳秒级 & 连续（锁定微秒级） & 硬件：鉴相器、VCO、延迟线 \\
DRAM 命令调度 & 总线空闲窗口，纳秒级 & 每时钟周期 & 硬件：内存控制器 \\
刷新调度 & 电容保持，毫秒级 & \hwevidence{7.8 $\mu$s}（高温 3.9 $\mu$s） & 硬件：内存控制器定时 \\
HDD 伺服 & 机械运动，kHz 级带宽 & 数十 kHz 采样 & 硬件+盘控固件 \\
TFC 飞行高度 & 加热器热惯性，毫秒级 & kHz 至 ms 级 & 盘控固件 \\
I/O 调度 & 请求到达间隔，微秒至毫秒级 & 微秒至毫秒级 & 软件：内核块层 \\
DVFS & 负载变化，毫秒至秒级 & 毫秒级 & 软件/固件：调频策略 \\
热管理、风扇 & 整机热惯性，秒级 & 百毫秒至秒级 & 软件/固件：EC、OS \\
\bottomrule
\end{longtable}

\begin{classicnote}
"采样频率远高于被控对象带宽"是离散控制的通用前提，与信号重建里的奈奎斯特定理在精神上相通：你要控制一个对象，就得比它变得快。工程上取十倍以上的频率裕量，才有空间放下控制律本身的延迟。把这条判据反过来读，就得到整机设计的一条分工线——比微秒快的一切环必须进硅片，比毫秒慢的一切环都可以写成代码，中间地带（kHz 级的机械环）属于固化在设备里的固件。
\end{classicnote}

这条边界也解释了前几节的安排：存储器件结构章第四节里 FR-FCFS 调度必须做成硬件状态机，因为 DRAM 总线的空闲窗口只有几纳秒，任何软件都赶不上；而同样是"调度"，块层的 I/O 调度器可以从容地用微秒级软件实现，因为它面对的对象（队列深度、请求到达）本身就慢三个量级。环放在哪里，是被控对象的速度决定的，不是实现者的偏好。

\section{七、整机控制地图}

本书从第二章的交叉耦合反相器开始，一路讲到的几乎每一个子系统，内部都至少有一个反馈环。现在把它们全部列在一张表里：环路的被控量、传感器、执行器、控制器类型、时间尺度，以及它在本书哪里讲过。这张表是本章的收尾，也是全书硬件部分的一张索引。

\begin{longtable}{L{0.13\linewidth}L{0.11\linewidth}L{0.13\linewidth}L{0.11\linewidth}L{0.12\linewidth}L{0.09\linewidth}L{0.17\linewidth}}
\toprule
环路 & 被控量 & 传感器 & 执行器 & 控制器类型 & 时间尺度 & 讲过的地方 \\
\midrule
VRM 电压环 & 核心电压 & 输出端分压采样 & 功率级占空比 & 模拟 PI / 滞环 & ns--$\mu$s & 电源时钟章 \\
PLL 相位环 & 时钟相位与频率 & 鉴相器 & VCO 控制电压 & 模拟锁相环 & ps--ns & 电源时钟章 \\
DLL 延迟环 & 时钟延迟 & 相位比较器 & 可变延迟线 & 数字延迟锁定 & ps--ns & 存储器件结构章第三节 \\
内存训练 & 采样点位置 & 参考图案读回比对 & 延迟线、驱动强度、ODT & 固件扫参状态机 & 一次性，ms & 存储器件结构章第三节 \\
DRAM 命令调度 & 总线利用率与延迟 & 请求队列、bank 状态 & 命令发射顺序 & FR-FCFS 硬件调度 & ns & 存储器件结构章第四节 \\
刷新调度 & 电容电荷保持 & 定时器、温度传感器 & REF 命令 & 硬件定时加温度补偿 & $\mu$s & 存储器件结构章第二节 \\
HDD 伺服 & 磁头位置（PES） & 伺服扇区 burst & 音圈电机电流 & 数字伺服、双级执行器 & kHz & 存储器件结构章第六节 \\
TFC 飞行高度 & 磁头--盘片间距 & 读信号幅度 & 滑块加热器功率 & 盘控固件闭环 & ms & 存储器件结构章第六节 \\
cache MSHR（未命中状态保持寄存器） & 在途 miss 数 & miss 计数、队列占用 & 合并、阻塞、回退 & 硬件队列管理 & ns & 主存与总线章、存储器件结构章第九节 \\
NVMe 队列流量控制 & 在途命令数 & 完成队列计数 & 提交节奏、门铃 & 指针差限深 & $\mu$s--ms & 主存与总线章、外设队列章 \\
OS 调度 tick & CPU 时间分配 & tick 中断 + 计时 & 上下文切换 & CFS & ms & 本章第六节 \\
DVFS & 性能与功耗 & 负载计数器、功耗计 & 频率与电压档位 & OS/固件调频策略 & ms--s & 本章第六节 \\
风扇与热管理 & 芯片温度 & 片内热敏二极管 & 风扇 PWM、降频 & 滞环 / PID & s & 本章、经典芯片与平台章 \\
\bottomrule
\end{longtable}

把这张表按时间尺度从上往下读，一个结构浮现出来：整台机器就是几十个控制环的嵌套，时间尺度从皮秒一直排到秒。快的环离物理近——VRM 贴着功率级，PLL 贴着时钟树，DRAM 调度贴着命令总线，它们必须做进硅片，因为被控对象的变化比任何软件的反应都快；慢的环离物理远——DVFS 看的是毫秒级的负载趋势，风扇看的是秒级的温度走向，调度器看的是毫秒级的公平性，它们可以放在固件和操作系统里，用可升级的代码实现。中间地带属于设备自己的固件：硬盘伺服和 TFC 运行在盘控处理器上，内存训练运行在开机固件里——它们比软件快、又不必像模拟环那样连续，恰好落在可编程但不可打扰的位置。

这张地图同时回答了本书开头隐含的那个问题：什么叫理解一台机器。一半答案是数据通路——数据从哪里来、经过哪些变换、到哪里去，从取指到执行，从内存层次到 DMA 回填；另一半答案是控制面——谁按什么规则让这一切发生：哪个环稳住电压，哪个环对齐相位，哪个环决定下一条命令，哪个环分配时间。数据通路决定机器能做什么，控制面决定机器在什么条件下持续地做下去。两者都看见了，机器才不再是黑箱。
